\chapter{Basic Text Analysis and
Statistics}\label{sec-basic-text-analysis-and-statistics}

\emph{Learning Objectives}

After reading this chapter, you should be able to:

\emph{Overview starts}

\begin{itemize}
\tightlist
\item
  Import, prepare, and clean text data in R.
\item
  Extract and analyze basic statistics for text, including word frequencies.
\item
  Represent text numerically using models like bag-of-words and TF-IDF.
\item
  Understand and apply the vector space model for comparing texts.
\item
  Initiate exploratory text analysis using visual and statistical tools.
\end{itemize}

\emph{Overview ends}

Computers work with numbers. They are incredibly efficient at this task,
which is what they were initially built for: mathematical
computations. Gradually, some of the real-world mechanics have been
adapted and shifted to digital use. When we think about typing text,
although there were once plenty of typewriters, we have now switched to
computers which allow us to lay text on a white canvas of pixels that
mimics a sheet of paper. While the vast majority of people write
text, just a tiny fraction analyze it, and the latter is our job here.
But in text analytics, text must be made ready for interpretation by computer, and we have the duty to ensure this process is done
properly. We act at a higher level and do not see the transformation take place,
but we must follow specific strategies and procedures to convert text
documents into vectors of numbers. This mapping between a sequence of
symbols and digits is called feature representation, feature extraction,
or feature encoding (Goldberg 2017).

Before we look at deeper meanings or more elaborate forms of analysis, there
is much basic information that can be extracted from text, such as the
number of words, the number of characters, the number of special
characters, the number of numbers, the average word length, and so on.
While all these features can be extracted, a widely used conceptual
model of text known as bag-of-words (BoW) focuses on the frequency of
words in a text. For applications of this concept to marketing, see,
for example, Netzer, Lemaire, and Herzenstein (2019) and Toubia et al. (2019).
The name bag implies that the words are simply thrown into a container with no
specific order, and this is a fitting description. Whenever we apply
this simple model, our attention is directed exclusively to the
occurrence of the words and their frequency, but not to the structure of
the text. More details on this model will be discussed below. Curiously enough, by dismantling the text's structure---removing
and modifying words or parts of them and eliminating punctuation marks---we can gain insights without having to read hundreds of pages simply by
analyzing a selection of its (unordered) words. Using the techniques
presented in the previous chapter, we will start by collecting a simple
text from the internet and transferring it to R, and then conducting some very
basic analysis on it. After that we will use the dataset mentioned in
Chapter~\ref{introduction-1}, apply initial reasoned text-processing steps, and transform it
into a matrix, enabling our computer to handle it efficiently.
Finally, we will analyze the outcome. What we are trying to achieve
is size reduction. We are aiming to filter
information and extract knowledge from huge numbers of documents. To accomplish this, we have to cut
and remove parts of the text, leaving only the essential (and
meaningful) elements.

\section{The basics: A short use
case}\label{the-basics---a-short-use-case}

Most of the time, we need to work with text downloaded from the
Internet. It is nearly impossible to find text that perfectly fits our
specific needs. Therefore, a text analyst first collects and cleans
the text, preparing it for its analysis.

Let us explore this through a simple but comprehensive example to demonstrate
the process. We will show how to extract text from the internet, save
it to our computer, and perform basic analysis with the commands
discussed in the previous chapter. The example uses a text on text
analytics (no pun intended) from Project Gutenberg. There are two
built-in functions in R for reading information from the internet if it is
not in tabular form: \texttt{readLines()} and \texttt{scan()}. We prefer
the latter for its greater flexibility.

We provide the URL of the text and store it in a variable, named
\texttt{etext}. The function parameters \texttt{what} and \texttt{sep}
are set to interpret the input as character strings (not numerical
values) and to split the text into separate vector cells for each
new line, identified by the metacharacter \texttt{\textbackslash{}n}.

\emph{Program code starts}

\begin{Shaded}
\begin{Highlighting}[]
\CommentTok{\# Set URL of the text file}
\NormalTok{url }\OtherTok{\textless{}{-}} \StringTok{"http://www.gutenberg.org/cache/epub/53/pg53.txt"}

\CommentTok{\# Read text as character strings, split by newline}
\NormalTok{etext }\OtherTok{\textless{}{-}} \FunctionTok{scan}\NormalTok{(url, }\AttributeTok{what =} \StringTok{"character"}\NormalTok{, }\AttributeTok{sep=}\StringTok{"}\SpecialCharTok{\textbackslash{}n}\StringTok{"}\NormalTok{)}
\end{Highlighting}
\end{Shaded}

\begin{verbatim}
Read 6855 lines.
\end{verbatim}

\emph{Program code ends}

The system indicates the number of `items' stored in the \texttt{etext}
variable, signifying that the text has been segmented into smaller
parts. These parts are then stored in a vector of 6,855 elements. Simply
typing \texttt{etext} into the console will display the entire imported
text, which might be extensive and impossible to read. To 
inspect the first six elements of \texttt{etext} briefly, we use the
\texttt{head()} function.

\emph{Program code starts}

\begin{Shaded}
\begin{Highlighting}[]
\FunctionTok{head}\NormalTok{(etext)}
\end{Highlighting}
\end{Shaded}

\begin{verbatim}
[1] "The Project Gutenberg eBook of Workshop on Electronic Texts:"            
[2] "Proceedings, 9-10 June 1992"                                             
[3] ""                                                                        
[4] "This ebook is for the use of anyone anywhere in the United States and"   
[5] "most other parts of the world at no cost and with almost no restrictions"
[6] "whatsoever. You may copy it, give it away or re-use it under the terms"  
[7] "of the Project Gutenberg License included with this ebook or online"     
\end{verbatim}

\emph{Program code ends}

So, this is the process by which the text was segmented into smaller
pieces. The function splits the text at each occurrence of the newline
metacharacter \textbackslash n, which, while invisible to our eyes,
is present in the \textbf{.txt} file's structure. Examining the \textbf{pg53.txt} file, we realize that not all parts
of this text are relevant for our analysis. In fact, certain sections
need to be excluded to focus only on the pertinent content.

It is apparent that the meaningful content begins with the word
`INTRODUCTION' and concludes with `END'. However, we must note that the
word `END' appears twice in the text. With the help of the
\texttt{grep()} function with the \texttt{value} parameter set to TRUE,
we can isolate the existence of lines containing these key words.

\emph{Program code starts}

\begin{Shaded}
\begin{Highlighting}[]
\CommentTok{\# Identify lines containing \textquotesingle{}INTRODUCTION\textquotesingle{} or \textquotesingle{}}\RegionMarkerTok{END}\CommentTok{\textquotesingle{} and store the content}
\NormalTok{keyword\_lines }\OtherTok{\textless{}{-}} \FunctionTok{grep}\NormalTok{(}\StringTok{"}\SpecialCharTok{\textbackslash{}\textbackslash{}}\StringTok{bINTRODUCTION}\SpecialCharTok{\textbackslash{}\textbackslash{}}\StringTok{b|}\SpecialCharTok{\textbackslash{}\textbackslash{}}\StringTok{bEND}\SpecialCharTok{\textbackslash{}\textbackslash{}}\StringTok{b"}\NormalTok{, etext, }\AttributeTok{value =} \ConstantTok{TRUE}\NormalTok{)}
\NormalTok{keyword\_lines}
\end{Highlighting}
\end{Shaded}

\begin{verbatim}
[1] "INTRODUCTION"                                                        
[2] "END"                                                                 
[3] "*** END OF THE PROJECT GUTENBERG EBOOK WORKSHOP ON ELECTRONIC TEXTS:"
[4] "PROCEEDINGS, 9-10 JUNE 1992 ***"                                     
\end{verbatim}

\emph{Program code ends}

Here we want to isolate the section between `INTRODUCTION' and
the first occurrence of `END'. We can do that This can be done by setting the
\texttt{value} attribute to FALSE, which is equivalent to the default
value (i.e.,~the attribute not appearing in the function) while we try to
detect the indexes in the vector where these two words appear.

\emph{Program code starts}

\begin{Shaded}
\begin{Highlighting}[]
\CommentTok{\# Find indices of lines containing \textquotesingle{}INTRODUCTION\textquotesingle{} or \textquotesingle{}}\RegionMarkerTok{END}\CommentTok{\textquotesingle{}}
\NormalTok{keyword\_indices }\OtherTok{\textless{}{-}} \FunctionTok{grep}\NormalTok{(}\StringTok{"}\SpecialCharTok{\textbackslash{}\textbackslash{}}\StringTok{bINTRODUCTION}\SpecialCharTok{\textbackslash{}\textbackslash{}}\StringTok{b|}\SpecialCharTok{\textbackslash{}\textbackslash{}}\StringTok{bEND}\SpecialCharTok{\textbackslash{}\textbackslash{}}\StringTok{b"}\NormalTok{, etext)}
\NormalTok{keyword\_indices}
\end{Highlighting}
\end{Shaded}

\begin{verbatim}
[1]  110 6549 6556
\end{verbatim}

\emph{Program code ends}

Think back to our chest of drawers analogy. Here we are interested in the drawers labeled 110, 6549 and 6556. Whatever is
present in the text before index 110 and after index 6549 has to be
discarded. This is done by selecting the elements in the vector. In R we
use square brackets \texttt{{[}\ {]}} for indexing and the colon symbol
\texttt{:} is used to denote a range, like this:
\texttt{{[}\ m\ :\ n\ {]}}. In plain English, keep everything from index
110 to 6549. We thus overwrite the \texttt{etext} variable with the
newly selected text. Let us look at the updated \texttt{etext} variable.

\emph{Program code starts}

\begin{Shaded}
\begin{Highlighting}[]
\CommentTok{\# Select the elements between keyword\_indices}
\NormalTok{etext }\OtherTok{\textless{}{-}}\NormalTok{ etext[keyword\_indices[}\DecValTok{1}\NormalTok{]}\SpecialCharTok{:}\NormalTok{keyword\_indices[}\DecValTok{2}\NormalTok{]]}
\FunctionTok{head}\NormalTok{(etext)}
\end{Highlighting}
\end{Shaded}

\begin{verbatim}
[1] "                              INTRODUCTION"                               
[2] "The Workshop on Electronic Texts (1) drew together representatives of"    
[3] "various projects and interest groups to compare ideas, beliefs,"          
[4] "experiences, and, in particular, methods of placing and presenting"       
[5] "historical textual materials in computerized form.  Most attendees gained"
[6] "much in insight and outlook from the event.  But the assembly did not"    
\end{verbatim}

\emph{Program code ends}

However, perhaps working with individual words, rather than sentence fragments,
would be more efficient. This can be achieved by using
\texttt{strsplit()} with the \texttt{\textbackslash{}\textbackslash{}W}
argument, which splits the text at non-word characters, isolating the
words. This is the basic concept of tokenization, that of dividing text into
constituent words.

\emph{Program code starts}

\begin{Shaded}
\begin{Highlighting}[]
\CommentTok{\# Isolate words}
\NormalTok{etext }\OtherTok{\textless{}{-}} \FunctionTok{strsplit}\NormalTok{(etext, }\StringTok{"}\SpecialCharTok{\textbackslash{}\textbackslash{}}\StringTok{W"}\NormalTok{)}
\FunctionTok{head}\NormalTok{(etext)}
\end{Highlighting}
\end{Shaded}

\begin{verbatim}
[[1]]
 [1] ""             ""             ""             ""             ""            
 [6] ""             ""             ""             ""             ""            
[11] ""             ""             ""             ""             ""            
[16] ""             ""             ""             ""             ""            
[21] ""             ""             ""             ""             ""            
[26] ""             ""             ""             ""             ""            
[31] "INTRODUCTION"

[[2]]
 [1] "The"             "Workshop"        "on"              "Electronic"     
 [5] "Texts"           ""                "1"               ""               
 [9] "drew"            "together"        "representatives" "of"             

[[3]]
 [1] "various"  "projects" "and"      "interest" "groups"   "to"      
 [7] "compare"  "ideas"    ""         "beliefs" 

[[4]]
 [1] "experiences" ""            "and"         ""            "in"         
 [6] "particular"  ""            "methods"     "of"          "placing"    
[11] "and"         "presenting" 

[[5]]
 [1] "historical"   "textual"      "materials"    "in"           "computerized"
 [6] "form"         ""             ""             "Most"         "attendees"   
[11] "gained"      

[[6]]
 [1] "much"     "in"       "insight"  "and"      "outlook"  "from"    
 [7] "the"      "event"    ""         ""         "But"      "the"     
[13] "assembly" "did"      "not"     
\end{verbatim}

\emph{Program code ends}

We notice that after applying the function we end up with a different
structure: a list. So each vector element includes a number of other elements. But,
if we altered the structure again, we would create a new vector with the
function \texttt{unlist()}. In this all the words would appear in a
sequential manner from the first to the last, but without punctuation
marks.

After applying \texttt{strsplit()}, the structure of our data changes
from a single vector into a list. In this list, each original element,
previously consisting of a line of text, is now a vector containing this
line of text broken into individual words. We have multiple vectors
stored inside a list. So, each `drawer' contains a vector of words,
similar to a drawer containing other drawers. To further transform this
structure into something more manageable, we need to flatten the list
and create a single chest of drawers with one word per drawer. We use
the \texttt{unlist()} function to transform our list into a single,
sequential vector with a word in each `drawer' but without any
punctuation marks.

\emph{Program code starts}

\begin{Shaded}
\begin{Highlighting}[]
\CommentTok{\# Vectorize}
\NormalTok{etext }\OtherTok{\textless{}{-}} \FunctionTok{unlist}\NormalTok{(etext)}
\FunctionTok{head}\NormalTok{(etext)}
\end{Highlighting}
\end{Shaded}

\begin{verbatim}
[1] "" "" "" "" "" ""
\end{verbatim}

\emph{Program code ends}

The first six elements might be blank because of the presence of white spaces. Our
focus is not on these, so we remove them. The
translation of the code into plain English is: From the variable
\texttt{etext}, extract the indexes (remember you access the content of a specific
cell using a number in squared brackets),of those
``drawers'' which are not empty. The exclamation mark \texttt{!} is used
for negation. The line of code selects all the ``drawers'' which are not
empty and overwrites \texttt{etext} with the filtered data.
Re-inspecting the variable containing the vector we will find that that
the blank spaces have been removed.

\emph{Program code starts}

\begin{Shaded}
\begin{Highlighting}[]
\CommentTok{\# Remove empty spaces}
\NormalTok{etext }\OtherTok{\textless{}{-}}\NormalTok{ etext[}\SpecialCharTok{!}\NormalTok{etext}\SpecialCharTok{==}\StringTok{""}\NormalTok{]}
\FunctionTok{head}\NormalTok{(etext)}
\end{Highlighting}
\end{Shaded}

\begin{verbatim}
[1] "INTRODUCTION" "The"          "Workshop"     "on"           "Electronic"  
[6] "Texts"       
\end{verbatim}

\emph{Program code ends}

To extract information from this text let us perform some basic
analysis. We need to find out the vocabulary used and the total number
of words. To achieve this, we use the \texttt{length()} function, which
gives us the total count of words. By applying the
\texttt{unique()} function, we can ensure that each word is counted only
once, hence we isolate the vocabulary. In simple terms, the vocabulary is
the set of words used in a given text.

\emph{Program code starts}

\begin{Shaded}
\begin{Highlighting}[]
\CommentTok{\# Check text length}
\FunctionTok{length}\NormalTok{(etext)}
\end{Highlighting}
\end{Shaded}

\begin{verbatim}
[1] 63301
\end{verbatim}

\begin{Shaded}
\begin{Highlighting}[]
\CommentTok{\# Check unique words number}
\FunctionTok{length}\NormalTok{(}\FunctionTok{unique}\NormalTok{(etext))}
\end{Highlighting}
\end{Shaded}

\begin{verbatim}
[1] 6725
\end{verbatim}

\emph{Program code ends}

The results show a total of 63,301 words spread over repeats and single
values, but the vocabulary itself is actually a lot smaller, consisting of 6,725
terms. Let us sort these terms alphabetically and store them in a
variable. Then we can peek at a segment of this variable (which, remember,
in this case is a vector) somewhere between positions 975 and 1000.
Again, imagine this vector as a chest of drawers, where each term
occupies its own drawer.

\emph{Program code starts}

\begin{Shaded}
\begin{Highlighting}[]
\CommentTok{\# Sort words in alphabetical order}
\NormalTok{s\_etext }\OtherTok{\textless{}{-}} \FunctionTok{sort}\NormalTok{(}\FunctionTok{unique}\NormalTok{(etext))}

\CommentTok{\# Take a glimpse into the vector}
\NormalTok{s\_etext[}\DecValTok{975}\SpecialCharTok{:}\DecValTok{1000}\NormalTok{]}
\end{Highlighting}
\end{Shaded}

\begin{verbatim}
 [1] "boldness"    "bond"        "Bonita"      "book"        "Book"       
 [6] "bookmark"    "books"       "Books"       "Boolean"     "bordering"  
[11] "Borders"     "bore"        "born"        "Borrowing"   "Boston"     
[16] "both"        "Both"        "bother"      "bottlenecks" "bottom"     
[21] "Boulevard"   "bound"       "Bound"       "boundaries"  "boundary"   
[26] "bounds"     
\end{verbatim}

\emph{Program code ends}

We can see that there are variations in words appearing in plural forms, written with
capital or with lowercase letters, or with different inflections. We
know that those words represent the same concept, no matter their
capitalization, but for a computer these are different, distinct words
with nothing in common.

To assess the popularity of a specific word we need to consider its frequency
of use: the more often a word appears, the more popular it is. For this,
we utilize the \texttt{table()} function, which counts
each word's occurrence and presents all them as if they were in an imaginary table.
To check which are the most frequently used words, we then apply the
\texttt{sort()} function. To view the top 25 most frequently used words, we can
use the \texttt{head()} function.

\emph{Program code starts}

\begin{Shaded}
\begin{Highlighting}[]
\CommentTok{\# Sort words based on frequency}
\NormalTok{sorted\_etext }\OtherTok{\textless{}{-}} \FunctionTok{sort}\NormalTok{(}\FunctionTok{table}\NormalTok{(etext), }\AttributeTok{decreasing =} \ConstantTok{TRUE}\NormalTok{)}
\FunctionTok{head}\NormalTok{(sorted\_etext, }\DecValTok{25}\NormalTok{)}
\end{Highlighting}
\end{Shaded}

\begin{verbatim}
etext
       the         of        and         to          a         in       that 
      3589       2447       1758       1752       1355       1204       1021 
        is        for         on       with         be         as         it 
       740        716        471        443        441        420        390 
       one        are        not         or        The       text         an 
       362        348        346        342        338        300        285 
     which         by          s electronic 
       280        269        265        261 
\end{verbatim}

\emph{Program code ends}

We notice that among the first words of interest are the nouns
\emph{text} and \emph{electronic}. This is intuitive, given that the
document discusses a workshop on electronic texts. Thus, even without
prior knowledge of the document's content, we can identify these as key
words. In any English-language text, certain words such as articles, pronouns, and prepositions occur frequently, composing about 25\% of
the text. These statistics (and more) can be verified with corpus-analysis tools such as Sketch Engine (Kilgarriff et al. 2014).
Additionally, we note that this roughly follows Zipf's empirical law (Ferrer i Cancho and
Solé 2001). This states that the
frequency of a word is inversely proportional to its rank: The most
frequent word appears twice as often as the second one, three times as
often as the third, and so on. The larger the corpus, the closer the
distribution of words aligns with Zipf's law. To  confirm
this visually, we have plotted the distribution of these words. Although not perfect, the
graph exhibits a tendency that aligns with Zipf's law. For plotting, we
use the \texttt{plot()} function with parameters like \texttt{xlab} for
the horizontal axis (x-axis), \texttt{ylab} for the vertical axis
(y-axis), and \texttt{las} for the rotation of labels on the x-axis. The
suffix \texttt{lab} in these parameters is short for label.

\emph{Program code starts}

\begin{Shaded}
\begin{Highlighting}[]
\FunctionTok{plot}\NormalTok{(sorted\_etext[}\DecValTok{1}\SpecialCharTok{:}\DecValTok{100}\NormalTok{],}
     \AttributeTok{type =} \StringTok{"b"}\NormalTok{,                }\CommentTok{\# Both points and lines}
     \AttributeTok{xlab =} \StringTok{"Rank of Words"}\NormalTok{,}
     \AttributeTok{ylab =} \StringTok{"Frequency of Words"}\NormalTok{,}
     \AttributeTok{las =} \DecValTok{2}\NormalTok{, }
     \AttributeTok{cex.axis =} \FloatTok{0.8}\NormalTok{,            }\CommentTok{\# Slightly reduce axis label size}
\NormalTok{     )}
\end{Highlighting}
\end{Shaded}

\begin{figure}[H]

{\centering \includegraphics{4.Chapter_files/figure-pdf/freqtermsplot-1.pdf}

}

\caption{Top 100 Words Frequency Distribution}

\end{figure}%

\emph{Program code ends}

Let us analyze what we have noticed so far. Several elements in the
document do not add value to our text analysis. Numbers, for example,
can be removed or converted into their literal counterparts (like turning
9 into \emph{nine}). The decision is yours. When we
transformed the text from a list to a vector
containing only the words, we can see the punctuation marks have been removed.
They are unnecessary for an analysis using the BoW model.
Similarly, empty spaces can be stripped away.

The most frequently occurring words, are so frequent that they often do not
add significant value to the understanding of the text, especially when it is only a
computer parsing the text. These words are technically known as stop
words and can be removed. Example of such words are: \emph{or},
\emph{and}, \emph{for}, and \emph{the}. Depending on your expertise, you can
also remove additional words that need to be left out. It is also
essential to standardize word capitalization to ensure consistency, so
they should appear uniformly. For example, words like \emph{Book},
\emph{BOOK}, and \emph{book} represent essentially the same term and
should typically be converted all to lowercase to unify their
appearance. Inflected forms should be reduced to a common root. We will
tackle this topic in more detail later. Carrying out properly what is often known as pre-processing the text
can significantly enhance the information we extract from it.

\section{Data preparation and normalization (text
pre-processing)}\label{sec-data-preparation-and-normalization-text-preprocessing}

To summarize, collecting text is the starting point of the whole process. Given that text must be structured and made ready for
computer analysis, it is essential to normalize, or pre-process, it
(Meyer, Hornik, and Feinerer 2008; Miner et al. 2012). Following
text correction using regular expressions (see
Chapter~\ref{textual-data}), the main preliminary operations are:
(i) corpus preparation, (ii) lowercase transformation, (iii) stop-word
removal, (iv) implicit-word removal, (v) punctuation removal, (vi)
number removal or number-to-text translation, and vice versa, (vii)
inflectional-form reduction, and, (viii) white-space contraction. We
demonstrate each of them below in the sequence in which they will be
applied.

\subsection{Corpus preparation}\label{corpus-preparation}

For out text mining procedures we will use the \texttt{tm} package: This is an abbreviation of \textbf{t}ext \textbf{m}ining. The functions in
this section refer to this, unless specified otherwise. The
\texttt{tm} package works with a structure called \texttt{Corpus} (note that the plural of corpus is corpora). This is an
abstract concept consisting of a collection of texts that are to be
analyzed. By default \texttt{tm} uses a volatile corpus, or
\texttt{VCorpus}, which is a structure held in memory in an R object
and which vanishes when destroyed.

In R, the function \texttt{options()} allows you to modify a number of default
values, one of them is \texttt{stringsAsFactors} in which the default is set
to \texttt{TRUE} and turns strings that appear more than once into
factors.\footnote{Factors are essentially what in statistics is called a
  categorical variable. If you encounter the same text value, R will
  automatically interpret it as a factor. For text analysis purposes, we
  want to treat text strings individually rather than as categories. In
  this case we need to set \texttt{stringsAsFactors} to FALSE}

In this example, we are working with the dataset detailed in the
introductory chapter and available for download from the book's website.
To import this dataset, which is in CSV format, we use the
\texttt{read.csv()} function. If you have not followed the setup
instructions in the preparatory chapter, you might have a different path
for where you store your dataset, so please check that everything is in
place to prevent errors. The dataset counts several columns, but we will
focus mainly on the one containing the reviews and store it
in a variable called \texttt{prodreviews}. The next step is to convert the
`Content' column to character type. This is important to ensure the consistency of the types of data
because sometimes R may not infer the data
types we need to work with correctly. The text for analysis can be sourced from a
directory, a vector treating each document as an individual item (in our
case, the variable \texttt{prodreviews}), or a data frame. Here we will
use \texttt{VectorSource}, but for more information see Feinerer and
Hornik (2019). When you have a vector where each element is a different
product review, as in our case, you can create a \texttt{Corpus}
object from this vector using \texttt{VectorSource}.

To get a glimpse of the first document in our corpus we can use the function
\texttt{inspect()}. The use of double squared brackets
\texttt{{[}{[}\ {]}{]}} indicates that there is more than one dimension,
similar to accessing elements in a list (R Core Team 2023).
\texttt{inspect(corpus{[}{[}1{]}{]})} will display the contents of the
first document in the corpus.

\begin{Shaded}
\begin{Highlighting}[]
\CommentTok{\# Load the tm library}
\FunctionTok{library}\NormalTok{(tm)}

\CommentTok{\# Set string handling to non{-}factor}
\FunctionTok{options}\NormalTok{(}\AttributeTok{stringsAsFactors =} \ConstantTok{FALSE}\NormalTok{)}

\CommentTok{\# Load the Product Reviews dataset}
\NormalTok{prodreviews }\OtherTok{\textless{}{-}} \FunctionTok{read.csv}\NormalTok{(}\StringTok{"data/Product Reviews.csv"}\NormalTok{)}

\CommentTok{\# Convert \textquotesingle{}Content\textquotesingle{} column to character type}
\NormalTok{content }\OtherTok{\textless{}{-}} \FunctionTok{as.character}\NormalTok{(prodreviews}\SpecialCharTok{$}\NormalTok{Content)}

\CommentTok{\# Create a volatile corpus}
\NormalTok{docs }\OtherTok{\textless{}{-}} \FunctionTok{VCorpus}\NormalTok{(}\FunctionTok{VectorSource}\NormalTok{(content))}

\CommentTok{\# Inspect the first document in the corpus}
\FunctionTok{inspect}\NormalTok{(docs[[}\DecValTok{1}\NormalTok{]])}
\end{Highlighting}
\end{Shaded}

\textless{}\textgreater{} Metadata: 7 Content: chars: 1384

I love my fire and highly recommend it to anyone who wants to watch
videos (netflix, hulu, amazon), read ebooks (purchased or from the local
library), surf the net and play games. I work in the tech field and I
LOVE apple entertainment products (I own many apple products and at work
I work with several). I am very thrilled with my fire (I LOVE IT TOO!)
because it works great as an entertainment product (and more affordable
than my apple products). I also think the fire is a great product
because of Amazons clout and support. Any tech product I buy needs to be
useable for years and updateable during those years. The amazon market
(still growing) is and will be getting better all the time. Why? Sheer
number of devices sold by Amazon equals a large potential population of
people who will buy software equals software programmers wanting to be a
part of that market. Also in IT the biggest frustration is having an
older device and not being able to update it or not being able to use
new software or technology because the manufacturer does not update
their devices (software) or provide backwards support. I am not saying
that Amazon will continue to provide support (that remains to be seen),
but the size, longevity, profitability and attitude of pro consumer of a
company are determining factors and compared to other companies I would
bet that Amazon will not disappoint.

\emph{Program code ends}

Inspecting the first review in our dataset, the output reveals that
it is a plain-text document. This output includes a metadata argument
which contains 7 additional elements of information, such as the author and
language. If you need extra information, use the function
\texttt{meta(docs{[}{[}1{]}{]})}. We also note that the text contains
1,384 characters, including white spaces.

Let us start off with the real text processing. The function that does
all the transformations is called \texttt{tm\_map()} and is a generic
function that maps a specific function to the elements of a corpus.

\subsection{Lowercase transformation}\label{lowercase-transformation}

For humans, the distinction between the words \emph{hotel} and
\emph{Hotel} is minimal and anyway, they both refer to lodgings. For computers this is different. Words have no semantic
meaning to them, are just a sequence of letters, and if two words
differ only through capitalization, then the two are distinct. Converting all words to lowercase ensures uniformity. This
enables the accurate counting of words. Converting to uppercase is also an
option but the text becomes less readable and, in online communication,
is interpreted as shouting. Remember, collecting accurate word-based
statistics is essential in the BoW model. The function that can
accomplish this is \texttt{tolower()} and is part of the R base program.
There is an important technical detail to note, while \texttt{tm\_map()}
is designed to work with corpus objects, using it with some standard R
functions directly can alter the structure of the content in the corpus.
To ensure that the content keeps its proper format, such as
\texttt{PlainTextDocument}, we use the \texttt{content\_transformer()}
wrapper function. By wrapping \texttt{tolower()} with
\texttt{content\_transformer()}, we make sure that the text is converted
to lowercase while maintaining the structure of the corpus.

\emph{Program code starts}

\begin{Shaded}
\begin{Highlighting}[]
\CommentTok{\# Transform all the characters to lower case}
\NormalTok{docs }\OtherTok{\textless{}{-}} \FunctionTok{tm\_map}\NormalTok{(docs, }\FunctionTok{content\_transformer}\NormalTok{(tolower))}

\CommentTok{\# Inspect the first document in the corpus}
\FunctionTok{inspect}\NormalTok{(docs[[}\DecValTok{1}\NormalTok{]])}
\end{Highlighting}
\end{Shaded}

\textless{}\textgreater{} Metadata: 7 Content: chars: 1384

i love my fire and highly recommend it to anyone who wants to watch
videos (netflix, hulu, amazon), read ebooks (purchased or from the local
library), surf the net and play games. i work in the tech field and i
love apple entertainment products (i own many apple products and at work
i work with several). i am very thrilled with my fire (i love it too!)
because it works great as an entertainment product (and more affordable
than my apple products). i also think the fire is a great product
because of amazons clout and support. any tech product i buy needs to be
useable for years and updateable during those years. the amazon market
(still growing) is and will be getting better all the time. why? sheer
number of devices sold by amazon equals a large potential population of
people who will buy software equals software programmers wanting to be a
part of that market. also in it the biggest frustration is having an
older device and not being able to update it or not being able to use
new software or technology because the manufacturer does not update
their devices (software) or provide backwards support. i am not saying
that amazon will continue to provide support (that remains to be seen),
but the size, longevity, profitability and attitude of pro consumer of a
company are determining factors and compared to other companies i would
bet that amazon will not disappoint.

\texttt{Program\ code\ ends}

Observe that all the words are now lowercase. While
this might not be the most reader-friendly format for us, it is
effective for computer processing, and, anyway, we do not need to read
the text.

\subsection{Stop-word removal}\label{stop-words-removal}

You should already be aware that the most frequent word in English
texts, is the article \emph{the}. Comparing the word frequencies of two different texts with
different topics, we would notice that the first
50 words would be likely to have similar frequencies. Such words that appear
very frequently but are also uninformative, we repeat this again,
are called stop words (Fox 1989). Once the \texttt{tm} library is loaded,
a complete list of the stop words can be found by typing
\texttt{stopwords()} in R's console. As you can notice, the list
includes contracted form such as \emph{isn't}. This tells us that we
need to remove the stop words \textbf{before} eliminating the punctuation. To
remove the stop words (and others if needed) from our text, we use the
\texttt{removeWords()} function.

\texttt{Program\ code\ starts}

\begin{Shaded}
\begin{Highlighting}[]
\CommentTok{\# remove the stopwords and specify the language}
\NormalTok{docs }\OtherTok{\textless{}{-}} \FunctionTok{tm\_map}\NormalTok{(docs, }\FunctionTok{content\_transformer}\NormalTok{(removeWords), }\FunctionTok{stopwords}\NormalTok{(}\StringTok{"english"}\NormalTok{))}

\CommentTok{\# Inspect the first document in the corpus}
\FunctionTok{inspect}\NormalTok{(docs[[}\DecValTok{1}\NormalTok{]])}
\end{Highlighting}
\end{Shaded}

\textless{}\textgreater{} Metadata: 7 Content: chars: 1071

love fire highly recommend anyone wants watch videos (netflix, hulu,
amazon), read ebooks (purchased local library), surf net play games.
work tech field love apple entertainment products ( many apple products
work work several). thrilled fire ( love !) works great entertainment
product ( affordable apple products). also think fire great product
amazons clout support. tech product buy needs useable years updateable
years. amazon market (still growing) will getting better time. ? sheer
number devices sold amazon equals large potential population people will
buy software equals software programmers wanting part market. also
biggest frustration older device able update able use new software
technology manufacturer update devices (software) provide backwards
support. saying amazon will continue provide support ( remains seen),
size, longevity, profitability attitude pro consumer company determining
factors compared companies bet amazon will disappoint.

\texttt{Program\ code\ ends}

We notice that the stop words have vanished from our
text, and the character count has reduced. Of course, the text becomes less
and less understandable for us to read.

\subsection{Removing other words}\label{own-dictionary-words-removal}

Sometimes we might need to remove other additional words. This process
is often accompanied by some trial and error, along with
informed decision-making. In a dataset where a city name appears
extensively, it becomes so prevalent in the corpus that we would like to
remove it because it will not contribute any additional insights. The
same \texttt{removeWords()} function can be used again. Suppose we
decide to remove the word \texttt{will} from our text.

\emph{Program code starts}

\begin{Shaded}
\begin{Highlighting}[]
\CommentTok{\# Remove specific words from the corpus}
\NormalTok{docs }\OtherTok{\textless{}{-}} \FunctionTok{tm\_map}\NormalTok{(docs, }\FunctionTok{content\_transformer}\NormalTok{(removeWords), }\FunctionTok{c}\NormalTok{(}\StringTok{"will"}\NormalTok{))}

\CommentTok{\# Inspect the first document in the corpus}
\FunctionTok{inspect}\NormalTok{(docs[[}\DecValTok{1}\NormalTok{]])}
\end{Highlighting}
\end{Shaded}

\textless{}\textgreater{} Metadata: 7 Content: chars: 1055

love fire highly recommend anyone wants watch videos (netflix, hulu,
amazon), read ebooks (purchased local library), surf net play games.
work tech field love apple entertainment products ( many apple products
work work several). thrilled fire ( love !) works great entertainment
product ( affordable apple products). also think fire great product
amazons clout support. tech product buy needs useable years updateable
years. amazon market (still growing) getting better time. ? sheer number
devices sold amazon equals large potential population people buy
software equals software programmers wanting part market. also biggest
frustration older device able update able use new software technology
manufacturer update devices (software) provide backwards support. saying
amazon continue provide support ( remains seen), size, longevity,
profitability attitude pro consumer company determining factors compared
companies bet amazon disappoint.

\emph{Program code ends}

We notice that strictly the word \emph{will} has been removed. It will
not be removed though when it appears in any other combination, meaning
as a prefix, suffix or a in a middle of a word. Also, you will observe
that the number of characters has decreased and that we have some more
empty spaces, exactly where the word \texttt{will} was present.

\textbf{Punctuation removal.}~Using BoW, which is centered on
word frequencies, the punctuation has to be removed. Because this model
overlooks semantic nuances, punctuation does not add value.
\texttt{removePunctuation()} is the function designed for this task.

\texttt{Program\ code\ starts}

\begin{Shaded}
\begin{Highlighting}[]
\CommentTok{\# Remove punctuation}
\NormalTok{docs }\OtherTok{\textless{}{-}} \FunctionTok{tm\_map}\NormalTok{(docs, removePunctuation)}

\CommentTok{\# Inspect the first document in the corpus}
\FunctionTok{inspect}\NormalTok{(docs[[}\DecValTok{1}\NormalTok{]])}
\end{Highlighting}
\end{Shaded}

\textless{}\textgreater{} Metadata: 7 Content: chars: 1021

love fire highly recommend anyone wants watch videos netflix hulu amazon
read ebooks purchased local library surf net play games work tech field
love apple entertainment products many apple products work work several
thrilled fire love works great entertainment product affordable apple
products also think fire great product amazons clout support tech
product buy needs useable years updateable years amazon market still
growing getting better time sheer number devices sold amazon equals
large potential population people buy software equals software
programmers wanting part market also biggest frustration older device
able update able use new software technology manufacturer update devices
software provide backwards support saying amazon continue provide
support remains seen size longevity profitability attitude pro consumer
company determining factors compared companies bet amazon disappoint

\texttt{Program\ code\ ends}

\subsection{Removing or transforming numbers}\label{numbers-removal-or-transformation}

Numbers may or may not be relevant to our analyses, but most of the time
they are not. Nevertheless, through practical experience you might find it
beneficial to transform digits into their word equivalents (for example, \texttt{3} into \texttt{three}) or to remove them completely. A
function can be easily created to map between numbers and text, if such
a transformation is necessary. The function \texttt{removeNumbers()}
will simply remove all numerical characters from the text.

\texttt{Program\ code\ starts}

\begin{Shaded}
\begin{Highlighting}[]
\CommentTok{\# Remove numbers}
\NormalTok{docs }\OtherTok{\textless{}{-}} \FunctionTok{tm\_map}\NormalTok{(docs, removeNumbers)}

\CommentTok{\# Inspect the first document in the corpus}
\FunctionTok{inspect}\NormalTok{(docs[[}\DecValTok{1}\NormalTok{]])}
\end{Highlighting}
\end{Shaded}

\textless{}\textgreater{} Metadata: 7 Content: chars: 1021

love fire highly recommend anyone wants watch videos netflix hulu amazon
read ebooks purchased local library surf net play games work tech field
love apple entertainment products many apple products work work several
thrilled fire love works great entertainment product affordable apple
products also think fire great product amazons clout support tech
product buy needs useable years updateable years amazon market still
growing getting better time sheer number devices sold amazon equals
large potential population people buy software equals software
programmers wanting part market also biggest frustration older device
able update able use new software technology manufacturer update devices
software provide backwards support saying amazon continue provide
support remains seen size longevity profitability attitude pro consumer
company determining factors compared companies bet amazon disappoint

\texttt{Program\ code\ ends}

\subsection{Reduction of inflectional
forms}\label{sec-reduction-of-inflectional-forms}

To reduce inflectional forms in a text, processes like stemming
(Lovins 1968) and lemmatization (Plisson et al. 2004; Toman, Tesar, and
Jezek 2006; Chai 2023) are commonly used. A stemming algorithm refers to
a rule-based heuristic computational process that aims to make
words uniform by reducing them to a common root (or a common stem), depending on whether
we remove the prefix or the suffix. For example, a stemmer
would reduce words such as \emph{marketing},
\emph{marketed}, or \emph{marketization} to the word \emph{market}. The goal is to achieve this correctly, hopefully most of the time. The \texttt{tm} package uses
the \texttt{SnowballC} (Bouchet-Valat and Bouchet-Valat 2020) package
for stemming which implements Martin Porter's version of the stemming
algorithm (1980) which focuses on suffix removal rather than
prefix removal. This version has become standard for stemming in English. Be aware that the algorithm does not reduce words to a
known stem word but to a neologism, for example, the words
\emph{replacement}, \emph{replace}, or \emph{replacing} might be reduced to
\emph{replac}. The function is \texttt{stemDocument()}.
Stemming is suitable for applications such as simple search engines, or,
as in our case, preliminary text analysis.

\texttt{Program\ code\ starts}

\begin{Shaded}
\begin{Highlighting}[]
\CommentTok{\# Load the SnoballC library for stemming}
\FunctionTok{library}\NormalTok{(SnowballC)}

\CommentTok{\# Stemming}
\NormalTok{docs }\OtherTok{\textless{}{-}} \FunctionTok{tm\_map}\NormalTok{(docs, }\FunctionTok{content\_transformer}\NormalTok{(stemDocument))}

\CommentTok{\# Inspect the first document in the corpus}
\FunctionTok{inspect}\NormalTok{(docs[[}\DecValTok{1}\NormalTok{]])}
\end{Highlighting}
\end{Shaded}

\textless{}\textgreater{} Metadata: 7 Content: chars: 788

love fire high recommend anyon want watch video netflix hulu amazon read
ebook purchas local librari surf net play game work tech field love appl
entertain product mani appl product work work sever thrill fire love
work great entertain product afford appl product also think fire great
product amazon clout support tech product buy need useabl year updat
year amazon market still grow get better time sheer number devic sold
amazon equal larg potenti popul peopl buy softwar equal softwar programm
want part market also biggest frustrat older devic abl updat abl use new
softwar technolog manufactur updat devic softwar provid backward support
say amazon continu provid support remain seen size longev profit attitud
pro consum compani determin factor compar compani bet amazon disappoint

\texttt{Program\ code\ ends}

Lemmatization is the dictionary- and linguistics-based process of the
reducing words to their dictionary form, known as the lemma, by removing
inflectional endings using vocabulary and morphological analysis. For
example, the words \emph{do}, \emph{did}, \emph{done} are all reduced to
the lemma \emph{do}. The package \texttt{tm} does not include a specific
lemmatization function, but we could use the \texttt{textstem} package
for this purpose. Since a corpus is composed of strings rather than
individual words, the \texttt{lemmatize\_strings()} function from
\texttt{textstem} can be adapted for use with \texttt{tm}. In order to work well, lemmatizers need to understand the structure of the phrase, which is important for semantic preservation. Using a BoW approach this
is disregarded, but still, lemmatization can be useful to reduce the
size of the vocabulary. For training purposes, if you have your
\texttt{textstem} library loaded, you can apply lemmatization to your
corpus with the following code:
\texttt{docs\ \textless{}-\ tm\_map(docs,\ content\_tranformer(lemmatize\_strings))}.
Remember, this function works best when the structure of sentences is
intact. Lemmatization is ideal in applications such as sentiment
analysis, natural language understanding, or machine learning models.

\subsection{White-space deletion}\label{white-space-deletion}

The goal here is to remove the extra white space from a text document to
condense instances of multiple white-space characters to a single blank
space. The function which accomplishes this task is
\texttt{stripWhitespace()}.

\texttt{Program\ code\ starts}

\begin{Shaded}
\begin{Highlighting}[]
\CommentTok{\# White space removal}
\NormalTok{docs }\OtherTok{\textless{}{-}} \FunctionTok{tm\_map}\NormalTok{(docs, stripWhitespace)}

\CommentTok{\# Inspect the first document in the corpus}
\FunctionTok{inspect}\NormalTok{(docs[[}\DecValTok{1}\NormalTok{]])}
\end{Highlighting}
\end{Shaded}

\textless{}\textgreater{} Metadata: 7 Content: chars: 788

love fire high recommend anyon want watch video netflix hulu amazon read
ebook purchas local librari surf net play game work tech field love appl
entertain product mani appl product work work sever thrill fire love
work great entertain product afford appl product also think fire great
product amazon clout support tech product buy need useabl year updat
year amazon market still grow get better time sheer number devic sold
amazon equal larg potenti popul peopl buy softwar equal softwar programm
want part market also biggest frustrat older devic abl updat abl use new
softwar technolog manufactur updat devic softwar provid backward support
say amazon continu provid support remain seen size longev profit attitud
pro consum compani determin factor compar compani bet amazon disappoint

\texttt{Program\ code\ ends}

After all these steps, we have reached a point where we need to
introduce text representation, a transformation that converts
words and phrases into numbers or vectors, enabling computers to process
textual data. This numerical representation is essential for various methods of text
analysis, such as topic modeling, sentiment analysis, and
machine learning algorithms applied to text data. We will analyze all of
these in the following chapters.

\section{Text representation}\label{text-representation}

Text representation is linked to a language model. But what exactly is a
model? It is a simplified and interpretative representation of reality which selectively emphasizes certain features,
but is not reality. For example, in geography a model could be
a map, for transportation a train schedule, in music the notation, in finance or marketing graphs and/or diagrams
showing hypothetical relationships between variables, such as ad spend vs. sales. Models can do more than just describe reality,
they can also predict or forecast phenomena (e.g., weather forecasts), or reveal patterns (e.g., seasonality in sales data). Mathematical models might show
predictions of weather patterns or the understanding of latent information such as the
seasonality of sales through a time-series graph. In the context of climate
change, models can provide projections based on current data trends,
such as \(\text{CO}_2\) emissions. Statistical models are used to
summarize data and reveal patterns. These models use observations to
extract and infer information, helping with tasks such as classifying or
predicting. In a statistical language model (or simply a language
model), when you have a set of documents (a corpus), the model estimates
the probability that the text in the documents was generated by that
model (Aggarwal 2018). It is important to note that models cannot perfectly
mirror reality. They are abstractions which come laden with limitations and
assumptions. For example, a map sacrifices details of terrain for simplicity, and a language model may miss nuanced contextual meanings.


A statistical language model also implies that there are mathematical operations,
so the text needs to be transformed into mathematical values, shifting
from unstructured to structured tabular formats. Like all statistical models, a language model (LM),
with a specified language such as English, is designed to assign
probabilities to phrases. More precisely, LMs are designed to
calculate the probability or ``the likelihood of a given word (or a
sequence of words) to follow a sequence of words'' Beumer (2020). A simple but
frequently used model to represent text is the vector
space model (VSM). A VSM
encodes words into numerical vectors, in a similar way to how a map shows territory. However, VSMs treat words as independent dimensions, which oversimplifies semantic relationships (e.g., synonyms are not inherently linked with each other). Let us see what we can work with: We have
\emph{terms}, which are basically words after they have been processed (e.g., "running" into "run"),
\emph{documents}, which can be phrases, paragraphs, articles, chapters, or
books, and \emph{corpora}, which are a collection of these documents.

\subsection{Vector space model}\label{sec-vsm-bow}

In a VSM, the mapping is structured and gives mathematical meaning
to the terms, documents, and the corpus. We know that the terms in a
document appear with a certain frequency, which we can call weight
and denote with a \(w\). A single document can be denoted by a \(d\).
So, a document \(i\) in a collection of documents \(D\) that contains
\(N\) unique terms (\(N\) is the vocabulary's dimension) can be
represented mathematically as a vector:
\[d_i=(w_1, w_2, ..., w_j, ..., w_N)\] where \(1 \leq j \leq N\) and
\(i \in 1...D\) (Jurafsky and Martin 2008)

For example, a simple document \(d_1\) which contains the words ``The more the
better'' can be represented by \(d_1=(1,1,2)\) corresponding to
(``better'', ``more'', ``the''). This is a vectorial representation of
the text. If you put all these vectors together, one beneath the other,
the final result is a matrix. In the rows you will have the documents in
the corpus. In the columns you will have the words in the vocabulary.
Some words might appear in some documents but not in others, some
might appear in all the documents and some might appear in only one. This is the distribution of the words over the documents. The
matrix is our main structure from which we extract features. Maybe the
simplest technique in this model is the Bag-of-Words (BoW), where the
representation of the words is 1-gram, meaning that we consider one
word at a time. There are also other representations based on n-grams,
where \(n\) can take any number \(\geq 2\). In the BoW approach, it is
assumed that the important unit is the word, so structure and order can
be discarded. As such, the semantics of a document are given by the
frequency of the words used, and the similitude with other documents is
based on the similitude of word frequencies (Services 2015). BoW
performs well at representing and comparing the contents of documents.

To all intents and purposes, after the normalization process, we have developed a BoW model.
Despite its name, hinting that it is baggy, we still need a structure for it to
do computational tasks. This structure has a grid form and is called the
document-term matrix (\texttt{dtm}) (Hofmann 1999). The \texttt{dtm}
is a matrix with two dimensions, (\emph{r,c}) whose rows,
\emph{r}, represent the documents (in our case each parcel of review content) and
the columns, \emph{c}, the terms in the vocabulary. Note that the
vocabulary contains the total number of words that appear in the whole corpus
arranged alphabetically. In the \texttt{dtm}, the frequency of a term's
occurrence in a document is represented by \texttt{n} in the respective
row and column. A value of 0 indicates the absence of a term. The value
\emph{n} is given by the different representations of the text and there
are three ways to show it: one-hot encoding, frequency vectors (or
count vectors), and TF-IDF vectors.

\subsubsection{One-hot encoding}\label{one-hot-encoding}

Let us take a simple example from our dataset. This type of encoding uses
the binary system, if a word is present in the vocabulary, it will be
represented by a \texttt{1}, its absence by a \texttt{0}. However, this is inefficient due to the large number of
\texttt{0}s. See figure Figure~\ref{fig-hotencoding} for a
representation of this.

\begin{figure}

\centering{

\includegraphics[width=0.9\textwidth,height=\textheight]{figure/chapter4/vectorspacehot2.png}

}

\caption{\label{fig-hotencoding}Document-Term Matrix - One Hot-Encoding
\newline Authors' own figure}

\end{figure}%

\subsubsection{Frequency vectors}\label{frequency-vectors}

A matrix is similar to a mathematical table. Therows show the
documents in the order they appear, and the columns have the words
listed in alphabetical order. See Table~\ref{tbl-tfdtm} for a
``mathematical'' representation of the documents we downloaded
previously. This matrix tracks how often each word appears in each
document.

\begin{longtable}[]{@{}
  >{\raggedright\arraybackslash}p{(\columnwidth - 18\tabcolsep) * \real{0.0882}}
  >{\raggedright\arraybackslash}p{(\columnwidth - 18\tabcolsep) * \real{0.0735}}
  >{\raggedright\arraybackslash}p{(\columnwidth - 18\tabcolsep) * \real{0.1324}}
  >{\raggedright\arraybackslash}p{(\columnwidth - 18\tabcolsep) * \real{0.0735}}
  >{\raggedright\arraybackslash}p{(\columnwidth - 18\tabcolsep) * \real{0.1176}}
  >{\raggedright\arraybackslash}p{(\columnwidth - 18\tabcolsep) * \real{0.1471}}
  >{\raggedright\arraybackslash}p{(\columnwidth - 18\tabcolsep) * \real{0.1029}}
  >{\raggedright\arraybackslash}p{(\columnwidth - 18\tabcolsep) * \real{0.0735}}
  >{\raggedright\arraybackslash}p{(\columnwidth - 18\tabcolsep) * \real{0.1176}}
  >{\raggedright\arraybackslash}p{(\columnwidth - 18\tabcolsep) * \real{0.0735}}@{}}
\caption{Document-Term Matrix}\label{tbl-tfdtm}\tabularnewline
\toprule\noalign{}
\begin{minipage}[b]{\linewidth}\raggedright
\end{minipage} & \begin{minipage}[b]{\linewidth}\raggedright
\ldots{}
\end{minipage} & \begin{minipage}[b]{\linewidth}\raggedright
ability
\end{minipage} & \begin{minipage}[b]{\linewidth}\raggedright
abl
\end{minipage} & \begin{minipage}[b]{\linewidth}\raggedright
access
\end{minipage} & \begin{minipage}[b]{\linewidth}\raggedright
accident
\end{minipage} & \begin{minipage}[b]{\linewidth}\raggedright
added
\end{minipage} & \begin{minipage}[b]{\linewidth}\raggedright
\ldots{}
\end{minipage} & \begin{minipage}[b]{\linewidth}\raggedright
comput
\end{minipage} & \begin{minipage}[b]{\linewidth}\raggedright
\ldots{}
\end{minipage} \\
\midrule\noalign{}
\endfirsthead
\toprule\noalign{}
\begin{minipage}[b]{\linewidth}\raggedright
\end{minipage} & \begin{minipage}[b]{\linewidth}\raggedright
\ldots{}
\end{minipage} & \begin{minipage}[b]{\linewidth}\raggedright
ability
\end{minipage} & \begin{minipage}[b]{\linewidth}\raggedright
abl
\end{minipage} & \begin{minipage}[b]{\linewidth}\raggedright
access
\end{minipage} & \begin{minipage}[b]{\linewidth}\raggedright
accident
\end{minipage} & \begin{minipage}[b]{\linewidth}\raggedright
added
\end{minipage} & \begin{minipage}[b]{\linewidth}\raggedright
\ldots{}
\end{minipage} & \begin{minipage}[b]{\linewidth}\raggedright
comput
\end{minipage} & \begin{minipage}[b]{\linewidth}\raggedright
\ldots{}
\end{minipage} \\
\midrule\noalign{}
\endhead
\bottomrule\noalign{}
\endlastfoot
Doc1 & \ldots{} & 0 & 2 & 0 & 0 & 0 & \ldots{} & 0 & \ldots{} \\
Doc2 & \ldots{} & 0 & 0 & 0 & 0 & 0 & \ldots{} & 0 & \ldots{} \\
Doc3 & \ldots{} & 0 & 0 & 0 & 1 & 0 & \ldots{} & 0 & \ldots{} \\
Doc4 & \ldots{} & 2 & 0 & 3 & 0 & 0 & \ldots{} & 0 & \ldots{} \\
Doc5 & \ldots{} & 0 & 0 & 0 & 0 & 0 & \ldots{} & 0 & \ldots{} \\
Doc6 & \ldots{} & 0 & 0 & 0 & 0 & 0 & \ldots{} & 3 & \ldots{} \\
\ldots{} & \ldots{} & \ldots{} & \ldots{} & \ldots{} & \ldots{} &
\ldots{} & \ldots{} & \ldots{} & \ldots{} \\
\end{longtable}

For a visual representation that illustrates the
relationship between documents schematically, their respective words, and the
document-term matrix (DTM), see Figure~\ref{fig-dtmfv}.

\begin{figure}

\centering{

\includegraphics[width=0.9\textwidth,height=\textheight]{figure/chapter4/dtm.png}

}

\caption{\label{fig-dtmfv}Document-Term Matrix Frequency Vectors}

\end{figure}%

\subsubsection{TF-IDF}\label{tf-idf}

A model that which acquired a good reputation is 
the TF-IDF measure. This is used in text analysis for document
comparison or for text classification, but its origins were rooted in
text information retrieval long before computers were even popular
(Jones 1972, 1973; Aizawa 2003). Most importantly, it compares
documents and shows which is more important based on specific words, in a very similar way to
how a search engine prioritizes information retrieval based on search
terms. The reason why TF-IDF works well with a BoW approach,

The basic use of TF-IDF is to assess the importance of a term in a
document. This is based on statistical principles. Remember, statistics
is the branch of mathematics that, among many other operations, also analyzes
the frequency of a (past) event. So the more important a term is, the
`heavier', or more frequent, it is expected to be.

To get an intuitive understanding, of how the algorithm works, note that it is composed of two
parts, the TF (\textbf{T}erm \textbf{f}requency) and the IDF
(\textbf{i}nverse \textbf{d}ocument \textbf{f}requency). TF quantifies
how frequently a term appears in a specific document relative to all
other words in that document: the more frequent a word is,
the higher the TF. Conversely, IDF instead the rarity of a word in
\textbf{a corpus of documents}. A word that appears in fewer documents
is more significant because it is rare, and thus has a higher IDF value.
The product of TF and IDF gives us a score (or a weight) for each word,
ranking its importance with respect to the other words. Let us look at an
example and examine the procedure in detail. Consider two sentences:

\begin{itemize}
\tightlist
\item
  Great tablet for the price.
\item
  The best tablet.
\end{itemize}

The sentences contain 5 and 3 terms respectively, excluding punctuation.
We organize them alphabetically because it facilitates tracking and
searching for specific terms. The collection of all the terms in these sentences
forms the vocabulary.

In Table~\ref{tbl-tfidf} we see the vector representation of the two
sentences. The header row contains the terms. The first row shows the
number of documents in which each term appears. The second and the third rows
contain the relative \textbf{T}erm \textbf{f}requency for the first and
the second documents respectively. The word term, \(t\), is the same
as the word, \(w\). The formula to calculate \textbf{TF} is
\(TF_{(t,d)}=\frac{f_{t,d}}{\sum_{t^{'}\in d}f_{t^{'},d}}\). This means
the term frequency of \emph{t} in document \emph{d} is the ratio of
the term's frequency \emph{f} in that document to the total number of
terms in the document. That total is the sum of all the terms
\(\sum_{t^{'}\in d}f_{t^{'},d}\) present in a document
\emph{d}. In our case, for the first sentence (which we consider a
document) there are a total of 7 words. The word ``the'' appears twice,
so its (relative) frequency, \textbf{TF}, is \(\frac{2}{7}\). The fourth
row represents the \textbf{i}nverse \textbf{d}ocument \textbf{f}requency
for each word in the vocabulary. Its simplified formula (simplified
because sometimes after the \(log()\) a 1 is added to avoid edge
cases in text processing) is \(IDF_{(t)}=log(\frac{N}{n_{t}})\), where
\(N\) represents the total number of the documents, in our case 2, and
\(n_{t}\) is the number of documents in which the term appears (see row
one). For example, the word ``the'' appears in both documents, so
\(n_{t}=2\). From algebra we know how to compute
\(log(\frac{2}{2})=log(1)=0\). The fifth and the sixth rows represent
the multiplication of each document's \textbf{TF} with the corresponding
\textbf{IDF}.

\begin{longtable}[]{@{}lllllll@{}}
\caption{TF-IDF Table}\label{tbl-tfidf}\tabularnewline
\toprule\noalign{}
& best & for & great & price & tablet & the \\
\midrule\noalign{}
\endfirsthead
\toprule\noalign{}
& best & for & great & price & tablet & the \\
\midrule\noalign{}
\endhead
\bottomrule\noalign{}
\endlastfoot
\(n_t\) & 1 & 1 & 1 & 1 & 2 & 2 \\
TF Doc 1 & 0 & 0.2 & 0.2 & 0.2 & 0.2 & 0.2 \\
TF Doc 2 & 0.34 & 0 & 0 & 0 & 0.34 & 0.34 \\
IDF & 0.3 & 0.3 & 0.3 & 0.3 & 0 & 0 \\
TF-IDF \(\times\) Doc 1 & 0 & 0.06 & 0.06 & 0.06 & 0 & 0 \\
TF-IDF \(\times\) Doc 2 & 0.102 & 0 & 0 & 0 & 0 & 0 \\
\end{longtable}

There are several ways you can calculate \textbf{TF-IDF} but the
principle remains the same: for the TF, the more frequent a term,
the more important it is, see Beel, Langer, and Gipp (2017). In IDF, the
rarer a term is in the corpus (i.e.,~among all the documents), the more
important it is. The word ``the'' will appear many times so it
will have a high TF because it is the most-used term in the English
language. But this word will be penalized in the IDF because it appears
in all the documents: the more it appears the lower the score, so it has
an inverse (document) frequency.

What TF-IDF accomplishes is to find the words that are important to a
document in a corpus. While TF relates to the first part of a BoW
approach, the TF-IDF metric will compute which words characterize the individual
documents. To sum up, a word that appears frequently in one or just a
few documents, will have a high \textbf{TF}\textbf{\(\times\)IDF} score.

The \texttt{dtm} containing the TF-IDF score for our main example about
the tablets is shown in Table~\ref{tbl-tfidf}.

\begin{longtable}[]{@{}
  >{\raggedright\arraybackslash}p{(\columnwidth - 16\tabcolsep) * \real{0.1045}}
  >{\raggedright\arraybackslash}p{(\columnwidth - 16\tabcolsep) * \real{0.0746}}
  >{\raggedright\arraybackslash}p{(\columnwidth - 16\tabcolsep) * \real{0.1343}}
  >{\raggedright\arraybackslash}p{(\columnwidth - 16\tabcolsep) * \real{0.1194}}
  >{\raggedright\arraybackslash}p{(\columnwidth - 16\tabcolsep) * \real{0.1194}}
  >{\raggedright\arraybackslash}p{(\columnwidth - 16\tabcolsep) * \real{0.1493}}
  >{\raggedright\arraybackslash}p{(\columnwidth - 16\tabcolsep) * \real{0.1045}}
  >{\raggedright\arraybackslash}p{(\columnwidth - 16\tabcolsep) * \real{0.0746}}
  >{\raggedright\arraybackslash}p{(\columnwidth - 16\tabcolsep) * \real{0.1194}}@{}}
\toprule\noalign{}
\begin{minipage}[b]{\linewidth}\raggedright
\end{minipage} & \begin{minipage}[b]{\linewidth}\raggedright
\ldots{}
\end{minipage} & \begin{minipage}[b]{\linewidth}\raggedright
ability
\end{minipage} & \begin{minipage}[b]{\linewidth}\raggedright
abl
\end{minipage} & \begin{minipage}[b]{\linewidth}\raggedright
access
\end{minipage} & \begin{minipage}[b]{\linewidth}\raggedright
accident
\end{minipage} & \begin{minipage}[b]{\linewidth}\raggedright
added
\end{minipage} & \begin{minipage}[b]{\linewidth}\raggedright
\ldots{}
\end{minipage} & \begin{minipage}[b]{\linewidth}\raggedright
comput
\end{minipage} \\
\midrule\noalign{}
\endhead
\bottomrule\noalign{}
\endlastfoot
Doc 1 & \ldots{} & 0 & 0.0476 & 0 & 0 & 0 & \ldots{} & 0 \\
Doc 2 & \ldots{} & 0 & 0 & 0 & 0 & 0 & \ldots{} & 0 \\
Doc 3 & \ldots{} & 0 & 0 & 0 & 0.0358 & 0 & \ldots{} & 0 \\
Doc 4 & \ldots{} & 0.0599 & 0 & 0.0625 & 0 & 0 & \ldots{} & 0 \\
Doc 5 & \ldots{} & 0 & 0 & 0 & 0 & 0 & \ldots{} & 0 \\
Doc 6 & \ldots{} & 0 & 0 & 0 & 0 & 0 & \ldots{} & 0.1069 \\
\ldots{} & \ldots{} & \ldots{} & \ldots{} & \ldots{} & \ldots{} &
\ldots{} & \ldots{} & \ldots{} \\
\end{longtable}

While TF-IDF can exist as a standalone algorithm, it can also be applied
to a BoW by replacing the word counts with their TF-IDF
values. The resulting matrix becomes a weighted version of BoW, where
each value indicates the importance of a term in a document relative to
the whole corpus. In simple terms, while BoW focuses solely on
frequency, TF-IDF balances frequency with importance. This is why the
initial stage in the BoW model is essential, because it allows the
TF-IDF matrix to emphasize words that are both informative and
contextually relevant.

\subsection{Size reduction}\label{size-reduction}

Now that we explored the purpose of preparing the text and how the VSM
quantifies words, an important further step before we move on to more
advanced techniques to analyze text is size reduction. By now it is
evident that the matrices that are created through a VSM can be huge,
depending on the number of documents and the breadth of vocabulary.
These matrices are often sparse. meaning that they contain a lot of
zero values due to the absence of most words in most documents.

In specific cases, infrequent words do not bring much insight,
and when this happens they can be removed. However, in other cases the words have to
remain because, as you could see in the TF-IDF model, they might
characterize a document or enrich a text. This might give us some
insights and we suggest not to dismiss the possibility of unearthing
knowledge from them. The words that appear only once in a corpus are
technically called \emph{hapax legomena}, or \emph{hapax} for short. The
longer the text, the higher the probability you will see an increase in
\emph{hapax} (Jockers and Thalken 2020a).

\subsubsection{Sparse-terms removal}\label{sparse-terms-removal}

Following data preparation and normalization (see
Section~\ref{sec-data-preparation-and-normalization-text-preprocessing}),
another procedure involves removing the sparse terms from
the \texttt{dtm}. These are terms that appear very infrequently in a
document and are recognized in a \texttt{dtm} by the fact that it
contains many entries set to zero in the \(j^{th}\) column and just a
few entries greater than zero. To remove them we can use
the sparse option from the \texttt{removeSparseTerms()} function. The
`sparsity' parameter here refers to a frequency threshold beyond which
terms are discarded. Sparsity decreases as it approaches 1.0 and can
hold values in the open interval (0, 1.0).

Let us take an example. If we look at the sample, notice the abundance of
zeros in the matrix. If the dataset was larger, words like ``ability''
or ``access'' would appear more frequently, but there would still be a lot of zeros in the columns. If a term does not have a lot of values
(again, in the columns), then the sparsity is increased, bringing the index
close to 1, meaning that there are a lot of zeros. If the word is
present in all the documents, for example the word ``tablet'', then
it will fill out its column with values different from zero, making the
sparsity index close to 0. We can conclude by saying that the sparsity
index essentially acts as a threshold, eliminating terms whose
proportion of zeros exceeds a specified level. For large corpora, a
threshold of between \(0.95\) and a value just below \(1\) is commonly
used.

\subsubsection{Word frequencies}\label{word-frequencies}

After the text has been prepared, a first simple descriptive analysis can be
carried out in order to search for words that most effectively represent
the content (Baayen 2001; Carroll and Roeloffs 1969). One key point is
to find the frequency for each term. We calculate this by adding up the occurences of a
term by column in a \texttt{dtm}. The terms are then
sorted in decreasing order based on their frequency. For an overview of
the data, we have shown the 20 most frequent terms along with their
corresponding frequency.

This simple form of text analysis can be visualized in the form of a
list, or in more visually appealing formats like plots or word clouds
(Heimerl et al. 2014). Frequently, co-occurring words can be associated
with each other (Church and Hanks 1990) and can be
represented graphically. This type of analysis gives us our first insights into the
patterns of word co-occurrence.

Here we see the code for the term-document matrix creation, the removal of
the sparse terms, and the extraction of the 20 most frequent words in
that corpus. To handle the sparse terms efficiently we use the library
\texttt{slam}. The function \texttt{row\_sums()} helps us to identify the
most frequent words.

\texttt{Program\ code\ starts}

\begin{Shaded}
\begin{Highlighting}[]
\FunctionTok{library}\NormalTok{(slam)}

\CommentTok{\# Create the term{-}document matrix}
\NormalTok{tdm }\OtherTok{\textless{}{-}} \FunctionTok{TermDocumentMatrix}\NormalTok{(docs)}

\CommentTok{\# Remove sparse terms with a threshold of 0.9}
\NormalTok{tdms }\OtherTok{\textless{}{-}} \FunctionTok{removeSparseTerms}\NormalTok{(tdm, }\FloatTok{0.9}\NormalTok{)}

\CommentTok{\# Sort the values in decreasing order and create a dataframe}
\NormalTok{v }\OtherTok{\textless{}{-}} \FunctionTok{sort}\NormalTok{(}\FunctionTok{row\_sums}\NormalTok{(tdms),}\AttributeTok{decreasing=}\ConstantTok{TRUE}\NormalTok{)}
\NormalTok{top\_terms }\OtherTok{\textless{}{-}} \FunctionTok{data.frame}\NormalTok{(}\AttributeTok{word =} \FunctionTok{names}\NormalTok{(v),}\AttributeTok{freq=}\NormalTok{v)}

\CommentTok{\# Display the top 20 terms}
\FunctionTok{head}\NormalTok{(top\_terms, }\DecValTok{20}\NormalTok{)}
\end{Highlighting}
\end{Shaded}

\begin{verbatim}
           word  freq
tablet   tablet 49184
use         use 46496
kindl     kindl 39523
fire       fire 36301
app         app 35770
can         can 30986
get         get 29015
ipad       ipad 28903
screen   screen 28536
one         one 27410
devic     devic 27333
like       like 27106
work       work 26595
just       just 23815
amazon   amazon 22473
great     great 22357
read       read 22115
android android 19169
want       want 18971
book       book 17854
\end{verbatim}

\texttt{Program\ code\ ends}

\subsubsection{Word bar plotting}\label{word-barplot}

To create a visual display of the most frequently employed words, a bar
plot can be used (Popescu, Mačutek, and Altmann 2009).
It is not feasible to plot all the words in the matrix due to the huge
size of the corpus. Instead, a frequency threshold is set to select
and display only the most prevalent terms. The threshold varies from
corpus to corpus, being linked to the volume of text and the most desirable number
of words. It is important to visualize the most frequent
words because they will be the most important words of the topics
treated in the text. The
words are shown in decreasing order on the x axis and their respective
occurrence depicted on the y axis, see
Figure~\ref{fig-most-frequent-words}. A bar plot allows us to identify
significant words immediately. Often, the most frequent words can point to
discussions about product features, and words like ``tablet'' and
``kindle'' identify key topics discussed in the corpus.

\texttt{Program\ code\ starts}

\begin{Shaded}
\begin{Highlighting}[]
\FunctionTok{barplot}\NormalTok{(}
\NormalTok{  top\_terms[}\DecValTok{1}\SpecialCharTok{:}\DecValTok{20}\NormalTok{, ]}\SpecialCharTok{$}\NormalTok{freq, }
  \AttributeTok{las =} \DecValTok{2}\NormalTok{,                         }
  \AttributeTok{names.arg =}\NormalTok{ top\_terms[}\DecValTok{1}\SpecialCharTok{:}\DecValTok{20}\NormalTok{, ]}\SpecialCharTok{$}\NormalTok{word,}
  \AttributeTok{col =} \StringTok{"lightblue"}\NormalTok{, }
  \AttributeTok{main =} \StringTok{"Most Frequent Words"}\NormalTok{,}
  \AttributeTok{ylab =} \StringTok{"Word Frequencies"}\NormalTok{,}
  \AttributeTok{cex.names =} \FloatTok{0.8}\NormalTok{,                 }
  \AttributeTok{cex.lab =} \FloatTok{0.9}\NormalTok{,                   }
  \AttributeTok{cex.main =} \DecValTok{1}                     
\NormalTok{)}
\end{Highlighting}
\end{Shaded}

\begin{figure}[H]

\centering{

\includegraphics{4.Chapter_files/figure-pdf/fig-most-frequent-words-1.pdf}

}

\caption{\label{fig-most-frequent-words}Most Frequent Words}

\end{figure}%

\texttt{Program\ code\ ends}

We can see that the nouns \emph{tablet}, \emph{kindl(e)},
\emph{fire}, \emph{screen}, etc. with their high frequency in the text
already provide significant clues to what the corpus refers to. These
words, being semantically related, point to what is predominant in these
documents, the Kindle Fire tablet.

\subsubsection{Word clouds}\label{wordcloud}

A word cloud is an effective tool for providing a quick and
direct visual overview of the most important words across a set of
documents. It is a popular method for summarizing qualitative data (W.
Cui et al. 2010; Scanfeld, Scanfeld, and Larson 2010). In marketing, word clouds can help identify key issues and
themes (Abell and Jones 2010), while in social-media analysis they offer
insights into prevailing topics and language use (Xu et al. 2013). In
the process of creating a word cloud, huge corpora of texts are reduced
to the words that occur frequently. There is a direct relationship
between the words and their appearance, bigger fonts are used for words
with a higher frequency, while smaller fonts depict a lower
frequency (see Figure~\ref{fig-wordclouds}). In such static
representations, words are considered as stand-alone entities without
defined associations or relationships. The plot is created using the \texttt{wordcloud} (Fellows et al. 2018) package. This
visualization underscores the importance of VSMs in text analysis.

\texttt{Program\ code\ starts}

\begin{Shaded}
\begin{Highlighting}[]
\FunctionTok{library}\NormalTok{(wordcloud)}

\CommentTok{\# Set seed for reproducibility}
\FunctionTok{set.seed}\NormalTok{(}\DecValTok{1234}\NormalTok{)}

\CommentTok{\# Create word cloud}
\FunctionTok{wordcloud}\NormalTok{(}
    \AttributeTok{words =}\NormalTok{ top\_terms}\SpecialCharTok{$}\NormalTok{word,    }\CommentTok{\# Words for the cloud}
    \AttributeTok{freq =}\NormalTok{ top\_terms}\SpecialCharTok{$}\NormalTok{freq,     }\CommentTok{\# Their frequencies}
    \AttributeTok{min.freq =} \DecValTok{1}\NormalTok{,              }\CommentTok{\# Minimum frequency for inclusion}
    \AttributeTok{max.words =} \DecValTok{200}\NormalTok{,           }\CommentTok{\# Limit on number of words}
    \AttributeTok{random.order =} \ConstantTok{FALSE}\NormalTok{,      }\CommentTok{\# Keep words in frequency order}
    \AttributeTok{rot.per =} \FloatTok{0.35}\NormalTok{,            }\CommentTok{\# Fraction of vertical words}
    \AttributeTok{scale =} \FunctionTok{c}\NormalTok{(}\FloatTok{3.5}\NormalTok{, }\FloatTok{0.25}\NormalTok{)       }\CommentTok{\# Scale for word size}
\NormalTok{)}
\end{Highlighting}
\end{Shaded}

\begin{figure}[H]

\centering{

\includegraphics{4.Chapter_files/figure-pdf/fig-wordclouds-1.pdf}

}

\caption{\label{fig-wordclouds}Word clouds}

\end{figure}%

\texttt{Program\ code\ ends}

We notice, as in the previous section, that words like
\emph{tablet}, \emph{kindl}, \emph{fire}, \emph{app}, \emph{screen},
\emph{amazon}, and \emph{work} are prominent. These frequently used
words highlight the focus of customer reviews in the dataset.
Adjectives such as \emph{great} hint at positive sentiment unless they are used
in a negative context. The nouns predominantly center around topics
related to the Fire tablet. From this small analysis, we can infer, even
without reading every review, that the Amazon Kindle tablet is a
generally well-received device. Colors in word clouds are used for
aesthetic purposes or to categorize words into different themes or
sentiments, not necessarily to indicate frequency. Word clouds, however,
are just the starting point for more comprehensive, in-depth analyses.

\subsubsection{Word correlation}\label{word-correlation}

Words in a text are related. They are linked to other words to create
meaning. To get deeper insights into our data, we can check for
correlations between some of the most frequent terms and the other words
in the document set (Kwartler 2017). By correlation, in this context, we
mean a measure of the words that co-occur in a document set. This
can indicate a relationship or association between the
terms. To analyze correlation you can use the \textbf{findAssocs()}
function from the \textbf{tm} package. The function requires you to specify
the document-term matrix and the correlation threshold which determines
the minimum strength of association to be considered significant, and is set to
a number between 0 and 1. If we set a correlation limit of \(0.10\), the
function will return words that co-occur with the target word with a
correlation score of \textbf{at least} \(0.10\). This number represents
the likelihood of the words appearing together relative to their occurring
independently. From an analytical point of view, using this we can
unveil potential thematic or contextual connections within the
documents. The correlation score computed by \texttt{findAssocs()} is
based on Pearson's correlation coefficient. This
evaluates the linear relationship between the frequency vector of the
target term and the frequency vectors of other terms in the
document-term matrix. This method highlights terms that exhibit a
significant linear relationship with the target word.

\texttt{Program\ code\ starts}

\begin{Shaded}
\begin{Highlighting}[]
\FunctionTok{library}\NormalTok{(tm)}
\FunctionTok{library}\NormalTok{(ggplot2)}
\FunctionTok{library}\NormalTok{(ggthemes)}

\CommentTok{\# Finding frequent terms}
\NormalTok{a }\OtherTok{\textless{}{-}} \FunctionTok{findFreqTerms}\NormalTok{(tdm, }\AttributeTok{lowfreq=}\DecValTok{10}\NormalTok{)}
\FunctionTok{head}\NormalTok{(a)  }\CommentTok{\# Display top terms}
\end{Highlighting}
\end{Shaded}

\begin{verbatim}
[1] "aaa"     "aac"     "aback"   "abandon" "abbrevi" "abc"    
\end{verbatim}

\begin{Shaded}
\begin{Highlighting}[]
\CommentTok{\# Finding word associations for "tablet", a correlation limit of 0.3}
\NormalTok{associations }\OtherTok{\textless{}{-}} \FunctionTok{findAssocs}\NormalTok{(tdm, }\AttributeTok{terms=}\StringTok{"tablet"}\NormalTok{, }\AttributeTok{corlimit=}\FloatTok{0.3}\NormalTok{)}
\CommentTok{\# Convert associations to dataframe}
\NormalTok{associations\_df }\OtherTok{\textless{}{-}} \FunctionTok{data.frame}\NormalTok{(}\AttributeTok{term =} \FunctionTok{names}\NormalTok{(associations}\SpecialCharTok{$}\NormalTok{tablet), }\AttributeTok{correlation =} \FunctionTok{unlist}\NormalTok{(associations}\SpecialCharTok{$}\NormalTok{tablet))}
\CommentTok{\# Sort associations by correlation in ascending order}
\NormalTok{associations\_df }\OtherTok{\textless{}{-}}\NormalTok{ associations\_df[}\FunctionTok{order}\NormalTok{(associations\_df}\SpecialCharTok{$}\NormalTok{correlation),]}


\CommentTok{\# Plotting associations}
\FunctionTok{ggplot}\NormalTok{(associations\_df, }\FunctionTok{aes}\NormalTok{(}\AttributeTok{x =}\NormalTok{ correlation, }\AttributeTok{y =} \FunctionTok{reorder}\NormalTok{(term, correlation))) }\SpecialCharTok{+}
  \FunctionTok{geom\_point}\NormalTok{(}\AttributeTok{color =} \StringTok{"blue"}\NormalTok{, }\AttributeTok{size =} \DecValTok{3}\NormalTok{) }\SpecialCharTok{+}
  \FunctionTok{geom\_segment}\NormalTok{(}\FunctionTok{aes}\NormalTok{(}\AttributeTok{x =} \DecValTok{0}\NormalTok{, }\AttributeTok{xend =}\NormalTok{ correlation, }\AttributeTok{y =}\NormalTok{ term, }\AttributeTok{yend =}\NormalTok{ term), }
               \AttributeTok{color =} \StringTok{"lightblue"}\NormalTok{, }\AttributeTok{size =} \FloatTok{0.7}\NormalTok{) }\SpecialCharTok{+}  \CommentTok{\# Add connecting lines for emphasis}
  \FunctionTok{theme\_minimal}\NormalTok{() }\SpecialCharTok{+}
  \FunctionTok{labs}\NormalTok{(}
    \AttributeTok{x =} \StringTok{"Correlation Score"}\NormalTok{,}
    \AttributeTok{y =} \StringTok{"Associated Terms"}
\NormalTok{  ) }\SpecialCharTok{+}
  \FunctionTok{theme}\NormalTok{(}
    \AttributeTok{axis.text =} \FunctionTok{element\_text}\NormalTok{(}\AttributeTok{size =} \DecValTok{10}\NormalTok{),}
    \AttributeTok{axis.title =} \FunctionTok{element\_text}\NormalTok{(}\AttributeTok{size =} \DecValTok{12}\NormalTok{),}
    \AttributeTok{plot.title =} \FunctionTok{element\_text}\NormalTok{(}\AttributeTok{size =} \DecValTok{14}\NormalTok{, }\AttributeTok{hjust =} \FloatTok{0.5}\NormalTok{),}
    \AttributeTok{panel.grid.major.y =} \FunctionTok{element\_blank}\NormalTok{()  }\CommentTok{\# Remove horizontal gridlines for a cleaner look}
\NormalTok{  )}
\end{Highlighting}
\end{Shaded}

\begin{figure}[H]

{\centering \includegraphics{4.Chapter_files/figure-pdf/assocciation-1.pdf}

}

\caption{Words association with `tablet'}

\end{figure}%

\texttt{Program\ code\ ends}

The results show connections between the word \emph{tablet} and the words
\emph{android}, \emph{use}, \emph{screen}, and so on. These findings
suggest that the terms frequently occur in the same context
giving more insight into what a user thinks about a product.

\subsubsection{Bi-gram word clouds}\label{bigram-wordclouds}

In a VSM, text can be easily tokenized as single words, having a
1-gram representation. We can also use 2-grams, commonly known as
bi-grams. You can have even longer combinations, called n-grams. A bi-gram
word cloud uses pairs of words as tokens, providing more insights into which terms are used together in reviews. A library that helps
tokenization is \textbf{RWeka} (Hornik, Buchta, and Zeileis 2009),
and allows for the creation of n-gram models.

\texttt{Program\ code\ starts}

\begin{Shaded}
\begin{Highlighting}[]
\FunctionTok{library}\NormalTok{(RWeka)}
\FunctionTok{library}\NormalTok{(wordcloud)}

\CommentTok{\# Define bigram tokenizer}
\NormalTok{Bigram\_Tokenizer }\OtherTok{\textless{}{-}} \ControlFlowTok{function}\NormalTok{(x) }\FunctionTok{NGramTokenizer}\NormalTok{(x, }\FunctionTok{Weka\_control}\NormalTok{(}\AttributeTok{min=}\DecValTok{2}\NormalTok{, }\AttributeTok{max=}\DecValTok{2}\NormalTok{))}

\CommentTok{\# Create bigram term{-}document matrix}
\NormalTok{bitdm }\OtherTok{\textless{}{-}} \FunctionTok{TermDocumentMatrix}\NormalTok{(docs, }\AttributeTok{control =} \FunctionTok{list}\NormalTok{(}\AttributeTok{tokenize =}\NormalTok{ Bigram\_Tokenizer))}

\CommentTok{\# Remove sparse terms}
\NormalTok{bitdms }\OtherTok{\textless{}{-}} \FunctionTok{removeSparseTerms}\NormalTok{(bitdm, }\FloatTok{0.99}\NormalTok{)}

\CommentTok{\# Convert to matrix and sort by frequency}
\NormalTok{v }\OtherTok{\textless{}{-}} \FunctionTok{sort}\NormalTok{(}\FunctionTok{rowSums}\NormalTok{(}\FunctionTok{as.matrix}\NormalTok{(bitdms)), }\AttributeTok{decreasing=}\ConstantTok{TRUE}\NormalTok{)}
\NormalTok{d }\OtherTok{\textless{}{-}} \FunctionTok{data.frame}\NormalTok{(}\AttributeTok{word =} \FunctionTok{names}\NormalTok{(v), }\AttributeTok{freq=}\NormalTok{v)}

\CommentTok{\# Generate bigram word cloud}
\FunctionTok{wordcloud}\NormalTok{(}\AttributeTok{words =}\NormalTok{ d}\SpecialCharTok{$}\NormalTok{word, }\AttributeTok{freq =}\NormalTok{ d}\SpecialCharTok{$}\NormalTok{freq, }\AttributeTok{random.order =} \ConstantTok{FALSE}\NormalTok{)}
\end{Highlighting}
\end{Shaded}

\begin{figure}[H]

{\centering \includegraphics{4.Chapter_files/figure-pdf/bigramswordcloud-1.pdf}

}

\caption{Word Bigrams}

\end{figure}%

\texttt{Program\ code\ ends}

It is evident that the words Kindle and Fire pair together, but other
words also provide insights: e.g., \emph{battery life}, \emph{touch
screen}, \emph{app store}, \emph{sd card}, and \emph{read books}.
Bi-grams offer more nuanced understanding than single words. The easiest
way to find how they are related to one another is through a word graph,
which illustrates how these terms are interconnected. An important
package that serves this purpose is \texttt{igraph} (Csardi and
Nepusz 2006). Using our existing dataframe containing
the bi-grams and their frequency, we can see that some words might also exist in other
combinations, and we need to separate them with the function
\texttt{separate()} from the \texttt{tidyr} (Wickham 2021) package. We can
filter out (see package \texttt{dplyr}; Wickham et al. 2021)) less
frequent combinations, such as those occurring fewer than 700 times, to give better readability. Adjusting this threshold can reveal different
patterns in the data. The combination of bi-grams will help you
understand that words in combination provide significant insights,
much more so than single words can provide. The power of word graphs is in the
visualization of how different words in a corpus are connected or
co-occur.

\begin{Shaded}
\begin{Highlighting}[]
\FunctionTok{library}\NormalTok{(dplyr)}
\FunctionTok{library}\NormalTok{(tidyr)}
\FunctionTok{library}\NormalTok{(igraph)}
\FunctionTok{library}\NormalTok{(ggraph)}

\CommentTok{\# Separate bigrams into two words}
\NormalTok{d\_2words}\OtherTok{\textless{}{-}}\NormalTok{ d }\SpecialCharTok{\%\textgreater{}\%}
  \FunctionTok{separate}\NormalTok{(word, }\AttributeTok{into =} \FunctionTok{c}\NormalTok{(}\StringTok{"word1"}\NormalTok{, }\StringTok{"word2"}\NormalTok{), }\AttributeTok{sep =} \StringTok{" "}\NormalTok{) }\SpecialCharTok{\%\textgreater{}\%}
  \FunctionTok{filter}\NormalTok{(freq }\SpecialCharTok{\textgreater{}} \DecValTok{700}\NormalTok{)}

\CommentTok{\# Build a filtered word network}
\NormalTok{word\_network }\OtherTok{\textless{}{-}} \FunctionTok{graph\_from\_data\_frame}\NormalTok{(d\_2words)}

\CommentTok{\# Inspect the word\_network variable}
\NormalTok{word\_network}
\end{Highlighting}
\end{Shaded}

\begin{verbatim}
IGRAPH e24d6a7 DN-- 110 96 -- 
+ attr: name (v/c), freq (e/n)
+ edges from e24d6a7 (vertex names):
 [1] kindl  ->fire      batteri->life      touch  ->screen    read   ->book     
 [5] sd     ->card      android->market    play   ->game      easi   ->use      
 [9] work   ->well      app    ->store     work   ->great     android->tablet   
[13] watch  ->movi      love   ->kindl     galaxi ->tab       power  ->button   
[17] can    ->get       custom ->servic    amazon ->prime     can    ->use      
[21] year   ->old       angri  ->bird      much   ->better    web    ->brows    
[25] android->app       usb    ->port      watch  ->video     high   ->recommend
[29] oper   ->system    even   ->though    surf   ->web       use    ->tablet   
+ ... omitted several edges
\end{verbatim}

The in the output is an internal identifier \texttt{igraph()} generates
to distinguish it within the R environment and is not essential for
the analysis. DN indicates a directed network (edges have directions).
There are 110 vertices (words) and 96 edges (word pairs). Each vertex
has an associated \texttt{name(v/c)} and each edge has an attribute
\texttt{freq\ (e/n)} indicating the frequency of the bi-gram.
The edges list displays word pairs.
\texttt{kindl\ -\textgreater{}\ fire} suggests a directed edge from
\texttt{kindl} to \texttt{fire}, these words frequently appearing
together in this order. The graph visualization will give a clearer
picture of the co-occurring words.

\begin{figure}[H]

{\centering \includegraphics{4.Chapter_files/figure-pdf/wordnetwork2-1.pdf}

}

\caption{Word Network}

\end{figure}%

The \texttt{ggraph()} function displays the information
contained in the word network beautifully. The legend shows the thickness of an
edge, representing the frequency of a bi-gram. Analyze all the words that
you see in the graph for a better understanding of the underlying
reviews. For example, we can notice words linked together such as
\emph{micro sd, sd card, card slot}, and credit card. The
direction of the arrow tells you which word comes first. You can also
see that Kindle Fire has the highest frequency bi-gram, and in the same
network segment these are linked to iPad. A lower threshold of the
filter frequency would give you different insights.

\subsubsection{Community Detection}\label{community-detection}

You can also attempt to perform community detection on the previous
graph. This is an important tool for uncovering patterns, clusters, or
hidden structures represented by words that have a higher probability of
appearing together in some groups than other groups. To
understand more about this type of analysis, see Fortunato and Hric
(2016). These communities often represent thematic or topical groupings
in texts, providing insights into the underlying structure and content of
the corpus. There are several community-detection methods, such as
Girvan-Newman, infomap, or label propagation (Lancichinetti and
Fortunato 2009) but a well-known one due to its fast parsing of large
datasets is the so-called Louvain method (Blondel et al. 2008). In word
networks, this is a popular community-detection tool used
to identify clusters of words that frequently co-occur in a corpus.
It iteratively groups nodes (words) into communities, so that the
modularity is maximized, revealing significant word clusters that might
represent specific topics or themes.

An interesting application of it for marketing is in market basket analysis (Faridizadeh, Abdolvand, and Rajaee Harandi
2018). This allows businesses to promote their
products more successfully and it is possible to discover communities with special
properties, gleaning knowledge from shopping behavior. In the
case of textual data, we obtain extra knowledge from how customers put
words together in their feedback and on which aspects they frequently focus.

\begin{Shaded}
\begin{Highlighting}[]
\CommentTok{\# Convert to undirected graph}
\NormalTok{word\_network }\OtherTok{\textless{}{-}} \FunctionTok{as.undirected}\NormalTok{(word\_network, }\AttributeTok{mode =} \StringTok{"collapse"}\NormalTok{)}

\CommentTok{\# Perform community detection using the Louvain method}
\NormalTok{communities }\OtherTok{\textless{}{-}} \FunctionTok{cluster\_louvain}\NormalTok{(word\_network, }\AttributeTok{resolution =} \DecValTok{1}\NormalTok{)}

\CommentTok{\# Visualization of the word network with communities}
\FunctionTok{ggraph}\NormalTok{(word\_network, }\AttributeTok{layout =} \StringTok{\textquotesingle{}fr\textquotesingle{}}\NormalTok{) }\SpecialCharTok{+}
  \FunctionTok{geom\_edge\_link}\NormalTok{(}\AttributeTok{edge\_alpha =} \FloatTok{0.5}\NormalTok{, }\AttributeTok{show.legend =} \ConstantTok{FALSE}\NormalTok{) }\SpecialCharTok{+}
  \FunctionTok{geom\_node\_point}\NormalTok{(}\FunctionTok{aes}\NormalTok{(}\AttributeTok{color =} \FunctionTok{factor}\NormalTok{(communities}\SpecialCharTok{$}\NormalTok{membership)), }\AttributeTok{size =} \DecValTok{3}\NormalTok{) }\SpecialCharTok{+}
  \FunctionTok{geom\_node\_text}\NormalTok{(}\FunctionTok{aes}\NormalTok{(}\AttributeTok{label =}\NormalTok{ name, }\AttributeTok{color =} \FunctionTok{factor}\NormalTok{(communities}\SpecialCharTok{$}\NormalTok{membership)),}
                 \AttributeTok{repel =} \ConstantTok{TRUE}\NormalTok{, }\AttributeTok{fontface =} \StringTok{"bold"}\NormalTok{, }\AttributeTok{size =} \DecValTok{2}\NormalTok{, }\AttributeTok{show.legend =} \ConstantTok{FALSE}\NormalTok{) }\SpecialCharTok{+}
  \FunctionTok{labs}\NormalTok{(}\AttributeTok{color =} \StringTok{"Community"}\NormalTok{) }\SpecialCharTok{+}
  \FunctionTok{theme\_void}\NormalTok{()}
\end{Highlighting}
\end{Shaded}

\begin{figure}[H]

{\centering \includegraphics{4.Chapter_files/figure-pdf/comunitydetection-1.pdf}

}

\caption{Word Network with Community Detection}

\end{figure}%

\subsubsection{Some considerations regarding word
frequencies}\label{some-considerations-on-word-frequencies}

The use of frequencies brings significant insights into the content of a
corpus, and word associations give a better understanding of their
co-occurrence in the text. Graphical representations are an easily grasped
schematic display of the word frequencies and their associations. While
these methods could be adopted for a preliminary analysis of the text,
they are not sufficient for discovering patterns, latent structures. or
groupings.

One of the most common activities of humans apears to be that of comparing and finding
similarities. Let us use this and go deeper into our analysis, making use of the word clouds in order
to compare two datasets (Kwartler 2017). As an exercise, we have split
our tablet dataset into reviews which are related to Kindle Fire (the
2011 version) and the 16GB iPad2, and have selected the first 100
reviews of both of them. We know that each product on Amazon is
identified by an ASIN number and we will use those IDs to focus on two
specific products. But before we do that, we have to build a function as a more
compact version of what we have been previously seen, something which does
the pre-processing directly on the imported documents. Previously we
acted on the corpus data with the help of the \texttt{tm\_map()}
function, now we will use the functions directly.

\texttt{Program\ code\ starts}

\begin{Shaded}
\begin{Highlighting}[]
\FunctionTok{library}\NormalTok{(tm)}
\FunctionTok{library}\NormalTok{(magrittr)}
\FunctionTok{library}\NormalTok{(wordcloud)}

\FunctionTok{options}\NormalTok{(}\AttributeTok{stringsAsFactors =} \ConstantTok{FALSE}\NormalTok{)}

\CommentTok{\# Load the data}
\NormalTok{prodreviews }\OtherTok{\textless{}{-}} \FunctionTok{read.csv}\NormalTok{(}\StringTok{"data/Product Reviews.csv"}\NormalTok{)}
\NormalTok{content }\OtherTok{\textless{}{-}} \FunctionTok{as.character}\NormalTok{(prodreviews}\SpecialCharTok{$}\NormalTok{Content)}

\CommentTok{\# Compact preprocessing (cleaning) function}
\NormalTok{preprocess }\OtherTok{\textless{}{-}} \ControlFlowTok{function}\NormalTok{(feedtext) \{}
\NormalTok{  feedtext }\SpecialCharTok{\%\textgreater{}\%}
    \FunctionTok{tolower}\NormalTok{() }\SpecialCharTok{\%\textgreater{}\%}
    \FunctionTok{removeWords}\NormalTok{(}\FunctionTok{stopwords}\NormalTok{(}\StringTok{"english"}\NormalTok{)) }\SpecialCharTok{\%\textgreater{}\%}
    \FunctionTok{removePunctuation}\NormalTok{() }\SpecialCharTok{\%\textgreater{}\%}
    \FunctionTok{stripWhitespace}\NormalTok{() }\SpecialCharTok{\%\textgreater{}\%}
    \FunctionTok{removeNumbers}\NormalTok{() }\SpecialCharTok{\%\textgreater{}\%}
    \FunctionTok{stemDocument}\NormalTok{() }\SpecialCharTok{\%\textgreater{}\%}
    \FunctionTok{return}\NormalTok{()}
\NormalTok{\}}

\CommentTok{\# Call the defined function}
\NormalTok{pre }\OtherTok{\textless{}{-}} \FunctionTok{preprocess}\NormalTok{(content)}

\CommentTok{\# Filter reviews by Item\_ID (ASIN) and select the first 100 comments for each product}
\NormalTok{kindle\_reviews }\OtherTok{\textless{}{-}} \FunctionTok{head}\NormalTok{(prodreviews}\SpecialCharTok{$}\NormalTok{Content[prodreviews}\SpecialCharTok{$}\NormalTok{Item\_ID }\SpecialCharTok{==} \StringTok{"B0051VVOB2"}\NormalTok{], }\DecValTok{100}\NormalTok{)  }\CommentTok{\# Kindle Fire 2011}
\NormalTok{ipad2\_reviews }\OtherTok{\textless{}{-}} \FunctionTok{head}\NormalTok{(prodreviews}\SpecialCharTok{$}\NormalTok{Content[prodreviews}\SpecialCharTok{$}\NormalTok{Item\_ID }\SpecialCharTok{==} \StringTok{"B0047DVWLW"}\NormalTok{], }\DecValTok{100}\NormalTok{)  }\CommentTok{\# iPad2 16GB}

\CommentTok{\# Preprocess the reviews}
\NormalTok{sub\_kindle }\OtherTok{\textless{}{-}} \FunctionTok{preprocess}\NormalTok{(}\FunctionTok{paste}\NormalTok{(kindle\_reviews, }\AttributeTok{collapse =} \StringTok{" "}\NormalTok{))}
\NormalTok{sub\_ipad2 }\OtherTok{\textless{}{-}} \FunctionTok{preprocess}\NormalTok{(}\FunctionTok{paste}\NormalTok{(ipad2\_reviews, }\AttributeTok{collapse =} \StringTok{" "}\NormalTok{))}

\CommentTok{\# Combine the two preprocessed review sets into a single vector}
\NormalTok{tablets }\OtherTok{\textless{}{-}} \FunctionTok{c}\NormalTok{(sub\_kindle, sub\_ipad2)}

\CommentTok{\# Create a corpus from the vector}
\NormalTok{corpus }\OtherTok{\textless{}{-}} \FunctionTok{VCorpus}\NormalTok{(}\FunctionTok{VectorSource}\NormalTok{(tablets))}

\CommentTok{\# Create a Term{-}Document Matrix}
\NormalTok{dtm }\OtherTok{\textless{}{-}} \FunctionTok{TermDocumentMatrix}\NormalTok{(corpus)}
\NormalTok{dtmmatrix }\OtherTok{\textless{}{-}} \FunctionTok{as.matrix}\NormalTok{(dtm)}

\CommentTok{\# Name the columns based on the products being compared}
\FunctionTok{colnames}\NormalTok{(dtmmatrix) }\OtherTok{\textless{}{-}} \FunctionTok{c}\NormalTok{(}\StringTok{"Kindle Fire 2011"}\NormalTok{, }\StringTok{"iPad2 16GB"}\NormalTok{)}
\end{Highlighting}
\end{Shaded}

\texttt{Program\ code\ ends}

Now the structure is ready for displaying the commonality cloud. The
word \emph{commonality} indicates that the words that customers use to
describe the two products are in common.

\texttt{Program\ code\ starts}

\begin{Shaded}
\begin{Highlighting}[]
\FunctionTok{commonality.cloud}\NormalTok{(dtmmatrix, }\AttributeTok{max.words=}\DecValTok{200}\NormalTok{, }\AttributeTok{random.order=}\ConstantTok{FALSE}\NormalTok{)}
\end{Highlighting}
\end{Shaded}

\begin{figure}[H]

{\centering \includegraphics{4.Chapter_files/figure-pdf/commonality-1.pdf}

}

\caption{Commonality Cloud}

\end{figure}%

\texttt{Program\ code\ ends}

We are trying to understand the linguistic patterns of two distinct sets of
documents related to the Kindle Fire and iPad2 and see what they have in common, how many
of these are discussed, and with which frequency. As bi-grams
are a common way to extract consecutive words, we have used the
\texttt{BigramTokenizer()} function from the \texttt{RWeka} package. The
custom function \texttt{extractBigrams()} is built to apply the
tokenizer across the corpus. The core of the analysis lies in comparing
the frequency of these common bi-grams. Be aware that 'lowercasing' and stemming are carried out during
pre-processing. The analysis
reveals the shared vocabulary between the two datasets and also the
relative emphasis or prevalence of certain themes or concepts,
shown by the frequency of the usage of the bi-grams.

\texttt{Program\ code\ starts}

\begin{Shaded}
\begin{Highlighting}[]
\FunctionTok{library}\NormalTok{(RWeka)}
\FunctionTok{library}\NormalTok{(dplyr)}
\FunctionTok{library}\NormalTok{(ggplot2)}
\FunctionTok{library}\NormalTok{(tm)}

\CommentTok{\# Bigram tokenizer function}
\NormalTok{BigramTokenizer }\OtherTok{\textless{}{-}} \ControlFlowTok{function}\NormalTok{(x) }\FunctionTok{NGramTokenizer}\NormalTok{(x, }\FunctionTok{Weka\_control}\NormalTok{(}\AttributeTok{min =} \DecValTok{2}\NormalTok{, }\AttributeTok{max =} \DecValTok{2}\NormalTok{))}

\CommentTok{\# Extract bigrams for a subset of text}
\NormalTok{extractBigrams }\OtherTok{\textless{}{-}} \ControlFlowTok{function}\NormalTok{(text) \{}
  \CommentTok{\# Extract bigrams and flatten the result to a single vector}
\NormalTok{  bigrams }\OtherTok{\textless{}{-}} \FunctionTok{unlist}\NormalTok{(}\FunctionTok{lapply}\NormalTok{(text, }\ControlFlowTok{function}\NormalTok{(line) }\FunctionTok{unlist}\NormalTok{(}\FunctionTok{BigramTokenizer}\NormalTok{(line))))}
  \FunctionTok{return}\NormalTok{(bigrams)}
\NormalTok{\}}

\CommentTok{\# Extract and process bigrams for Kindle Fire and iPad2}
\NormalTok{kindle\_bigrams }\OtherTok{\textless{}{-}} \FunctionTok{extractBigrams}\NormalTok{(sub\_kindle)}
\NormalTok{ipad2\_bigrams }\OtherTok{\textless{}{-}} \FunctionTok{extractBigrams}\NormalTok{(sub\_ipad2)}

\CommentTok{\# Calculate frequencies}
\NormalTok{kindle\_freq }\OtherTok{\textless{}{-}} \FunctionTok{sort}\NormalTok{(}\FunctionTok{table}\NormalTok{(kindle\_bigrams), }\AttributeTok{decreasing =} \ConstantTok{TRUE}\NormalTok{)}
\NormalTok{ipad2\_freq }\OtherTok{\textless{}{-}} \FunctionTok{sort}\NormalTok{(}\FunctionTok{table}\NormalTok{(ipad2\_bigrams), }\AttributeTok{decreasing =} \ConstantTok{TRUE}\NormalTok{)}

\CommentTok{\# Find common bigrams}
\NormalTok{common\_bigrams }\OtherTok{\textless{}{-}} \FunctionTok{intersect}\NormalTok{(}\FunctionTok{names}\NormalTok{(kindle\_freq), }\FunctionTok{names}\NormalTok{(ipad2\_freq))}

\CommentTok{\# Select top 20 common bigrams}
\NormalTok{top\_common\_bigrams }\OtherTok{\textless{}{-}} \FunctionTok{head}\NormalTok{(common\_bigrams, }\DecValTok{20}\NormalTok{)}

\CommentTok{\# Prepare data frame for plotting}
\NormalTok{bigram\_df }\OtherTok{\textless{}{-}} \FunctionTok{data.frame}\NormalTok{(}
  \AttributeTok{Bigram =} \FunctionTok{rep}\NormalTok{(top\_common\_bigrams, }\DecValTok{2}\NormalTok{),}
  \AttributeTok{Frequency =} \FunctionTok{c}\NormalTok{(}\FunctionTok{as.numeric}\NormalTok{(kindle\_freq[top\_common\_bigrams]), }
                \SpecialCharTok{{-}}\FunctionTok{as.numeric}\NormalTok{(ipad2\_freq[top\_common\_bigrams])), }\CommentTok{\# Make iPad2 frequencies negative}
  \AttributeTok{Product =} \FunctionTok{rep}\NormalTok{(}\FunctionTok{c}\NormalTok{(}\StringTok{"Kindle Fire 2011"}\NormalTok{, }\StringTok{"iPad2 16GB"}\NormalTok{), }\AttributeTok{each =} \FunctionTok{length}\NormalTok{(top\_common\_bigrams))}
\NormalTok{)}

\CommentTok{\# Create mirrored pyramid plot}
\FunctionTok{ggplot}\NormalTok{(bigram\_df, }\FunctionTok{aes}\NormalTok{(}\AttributeTok{x =} \FunctionTok{reorder}\NormalTok{(Bigram, }\SpecialCharTok{{-}}\FunctionTok{abs}\NormalTok{(Frequency)), }\AttributeTok{y =}\NormalTok{ Frequency, }\AttributeTok{fill =}\NormalTok{ Product)) }\SpecialCharTok{+}
  \FunctionTok{geom\_bar}\NormalTok{(}\AttributeTok{stat =} \StringTok{"identity"}\NormalTok{) }\SpecialCharTok{+}
  \FunctionTok{coord\_flip}\NormalTok{() }\SpecialCharTok{+}
  \FunctionTok{labs}\NormalTok{(}\AttributeTok{title =} \StringTok{"Comparison of Common Bigrams: Kindle Fire 2011 vs. iPad2 16GB"}\NormalTok{,}
       \AttributeTok{x =} \StringTok{"Bigram"}\NormalTok{,}
       \AttributeTok{y =} \StringTok{"Frequency"}\NormalTok{) }\SpecialCharTok{+}
  \FunctionTok{scale\_fill\_brewer}\NormalTok{(}\AttributeTok{palette =} \StringTok{"Set1"}\NormalTok{) }\SpecialCharTok{+}
  \FunctionTok{theme\_minimal}\NormalTok{() }\SpecialCharTok{+}
  \FunctionTok{theme}\NormalTok{(}\AttributeTok{panel.grid.major =} \FunctionTok{element\_blank}\NormalTok{(), }\AttributeTok{panel.grid.minor =} \FunctionTok{element\_blank}\NormalTok{(),}
        \AttributeTok{axis.text.y =} \FunctionTok{element\_text}\NormalTok{(}\AttributeTok{face =} \StringTok{"bold"}\NormalTok{))}
\end{Highlighting}
\end{Shaded}

\begin{figure}[H]

{\centering \includegraphics{4.Chapter_files/figure-pdf/bigramspyramid-1.pdf}

}

\caption{Top 20 Common Bi-grams for Kindle Fire2011 and iPad2 16GB}

\end{figure}%

\texttt{Program\ code\ ends}

In contrast, a comparison cloud makes a distinction between the datasets and shows what is specific to each. It takes two or more datasets and displays the frequency of the
words which are specific to each of them.

\texttt{Program\ code\ starts}

\begin{Shaded}
\begin{Highlighting}[]
\FunctionTok{comparison.cloud}\NormalTok{(}
\NormalTok{  dtmmatrix,}
  \AttributeTok{max.words =} \DecValTok{200}\NormalTok{, }
  \AttributeTok{random.order =} \ConstantTok{FALSE}\NormalTok{, }
  \AttributeTok{scale =} \FunctionTok{c}\NormalTok{(}\DecValTok{3}\NormalTok{, }\FloatTok{0.5}\NormalTok{),  }\CommentTok{\# Adjust the size scale for words}
  \AttributeTok{fixed.asp =} \ConstantTok{TRUE}    \CommentTok{\# Ensure the aspect ratio is fixed to make the cloud appear more square{-}like}
\NormalTok{)}
\end{Highlighting}
\end{Shaded}

\begin{figure}[H]

{\centering \includegraphics{4.Chapter_files/figure-pdf/comparison-1.pdf}

}

\caption{Comparison Cloud}

\end{figure}%

\texttt{Program\ code\ ends}

\subsubsection{Bi-grams comparison}\label{bigrams-comparison}

Another useful method can be the comparison of mutual bi-grams. With this
technique, you can acquire a quick overview of how bi-grams are used and
with what frequency, and an immediate comparison is possible between the two, see
Figure~\ref{fig-bigrams_comparison}. The steps for plotting the bi-grams are
as follows. First, extract them for both the Kindle Fire and iPad2 from the
pre-processed text. Then, identify the unique bigrams, i.e.,
the ones which are not in common to both datasets, by selecting the first
10. Eventually, put the selected data in a dataframe and plot them.
Reuse the function \texttt{extractBigrams()} from the previous
example.

\texttt{Program\ code\ starts}

\begin{Shaded}
\begin{Highlighting}[]
\CommentTok{\# Calculate frequencies}
\NormalTok{kindle\_freq }\OtherTok{\textless{}{-}} \FunctionTok{sort}\NormalTok{(}\FunctionTok{table}\NormalTok{(kindle\_bigrams), }\AttributeTok{decreasing =} \ConstantTok{TRUE}\NormalTok{)}
\NormalTok{ipad2\_freq }\OtherTok{\textless{}{-}} \FunctionTok{sort}\NormalTok{(}\FunctionTok{table}\NormalTok{(ipad2\_bigrams), }\AttributeTok{decreasing =} \ConstantTok{TRUE}\NormalTok{)}

\CommentTok{\# Identify unique bigrams}
\NormalTok{unique\_kindle }\OtherTok{\textless{}{-}} \FunctionTok{names}\NormalTok{(kindle\_freq)[}\SpecialCharTok{!}\FunctionTok{names}\NormalTok{(kindle\_freq) }\SpecialCharTok{\%in\%} \FunctionTok{names}\NormalTok{(ipad2\_freq)]}
\NormalTok{unique\_ipad }\OtherTok{\textless{}{-}} \FunctionTok{names}\NormalTok{(ipad2\_freq)[}\SpecialCharTok{!}\FunctionTok{names}\NormalTok{(ipad2\_freq) }\SpecialCharTok{\%in\%} \FunctionTok{names}\NormalTok{(kindle\_freq)]}

\CommentTok{\# Select top 10 unique bigrams}
\NormalTok{top\_kindle }\OtherTok{\textless{}{-}} \FunctionTok{head}\NormalTok{(}\FunctionTok{sort}\NormalTok{(kindle\_freq[unique\_kindle], }\AttributeTok{decreasing =} \ConstantTok{TRUE}\NormalTok{), }\DecValTok{10}\NormalTok{)}
\NormalTok{top\_ipad2 }\OtherTok{\textless{}{-}} \FunctionTok{head}\NormalTok{(}\FunctionTok{sort}\NormalTok{(ipad2\_freq[unique\_ipad], }\AttributeTok{decreasing =} \ConstantTok{TRUE}\NormalTok{), }\DecValTok{10}\NormalTok{)}

\CommentTok{\# Convert to data frames}
\NormalTok{df\_kindle }\OtherTok{\textless{}{-}} \FunctionTok{data.frame}\NormalTok{(}\AttributeTok{Bigram =} \FunctionTok{names}\NormalTok{(top\_kindle), }\AttributeTok{Frequency =} \FunctionTok{as.numeric}\NormalTok{(top\_kindle), }\AttributeTok{Product =} \StringTok{"Kindle Fire 2011"}\NormalTok{)}
\NormalTok{df\_ipad2 }\OtherTok{\textless{}{-}} \FunctionTok{data.frame}\NormalTok{(}\AttributeTok{Bigram =} \FunctionTok{names}\NormalTok{(top\_ipad2), }\AttributeTok{Frequency =} \FunctionTok{as.numeric}\NormalTok{(top\_ipad2), }\AttributeTok{Product =} \StringTok{"iPad2 16GB"}\NormalTok{)}

\CommentTok{\# Combine data frames}
\NormalTok{combined\_df }\OtherTok{\textless{}{-}} \FunctionTok{rbind}\NormalTok{(df\_kindle, df\_ipad2)}

\CommentTok{\# Plot using ggplot}
\FunctionTok{ggplot}\NormalTok{(combined\_df, }\FunctionTok{aes}\NormalTok{(}\AttributeTok{x =} \FunctionTok{reorder}\NormalTok{(Bigram, Frequency), }\AttributeTok{y =}\NormalTok{ Frequency, }\AttributeTok{fill =}\NormalTok{ Product)) }\SpecialCharTok{+}
  \FunctionTok{geom\_bar}\NormalTok{(}\AttributeTok{stat =} \StringTok{"identity"}\NormalTok{, }\AttributeTok{position =} \FunctionTok{position\_dodge}\NormalTok{()) }\SpecialCharTok{+}
  \FunctionTok{labs}\NormalTok{(}\AttributeTok{x =} \StringTok{"Bigram"}\NormalTok{, }\AttributeTok{y =} \StringTok{"Frequency"}\NormalTok{) }\SpecialCharTok{+}
  \FunctionTok{scale\_fill\_brewer}\NormalTok{(}\AttributeTok{palette =} \StringTok{"Set1"}\NormalTok{) }\SpecialCharTok{+}
  \FunctionTok{theme\_minimal}\NormalTok{() }\SpecialCharTok{+}
  \FunctionTok{coord\_flip}\NormalTok{()}
\end{Highlighting}
\end{Shaded}

\begin{figure}[H]

\centering{

\includegraphics{4.Chapter_files/figure-pdf/fig-bigrams_comparison-1.pdf}

}

\caption{\label{fig-bigrams_comparison}Top 10 Unique Bi-grams for Kindle
Fire 2011 and iPad2 16GB}

\end{figure}%

\texttt{Program\ code\ ends}

All these tools are important for understanding which features both products
have, or do not have, in common, how people use the words in
order to describe the products, and what makes them distinct. And it is not hard to imagine that
you might want to compare the reviews of what you produce or offer as a
service with the reviews of your competitors.

\subsection{Marketing applications}\label{marketing-applications}

How does the frequency of the words give insights into a set of
documents and how can this be used for marketing purposes? The frequency
of the words indicates that there is a tendency to describe a
phenomenon, following the argument that the more you talk about something, the more important it is. But
those words are perhaps the ones that are also used in a search
engine, so employing them, or at least considering them as SEO keywords, is as
important as checking trends with tools such as Google Trends. You
can be found on the internet if people look for you, so these words can
work as hooks. But how many words you should choose for SEO? Knowing
that the competition for one or two words is tight (see
Figure~\ref{fig-longtail}), the long tail (Anderson 2013) is a very
helpful place to examine. Check the search-engine results page and
you will see that Google's recent advances in NLP allow you to use a more generous number of words.

\begin{figure}

\centering{

\includegraphics[width=0.9\textwidth,height=\textheight]{4.Chapter_files/figure-pdf/fig-longtail-1.pdf}

}

\caption{\label{fig-longtail}The Long Tail for SEO Keywords}

\end{figure}%

By checking what you have in common with another company might help you diversify so that you can make
your distinctive. It is also important to analyze what already differentiates you from other companies, especially the keywords that put your products
in a bad light or if your products are compared to competitors'
products. You also need, in this case, to analyze not only the words
themselves but also associations of words. Another element to consider is
the time of the year, when is it better to market a product? Is it
perceived as a Christmas gift or a gadget for Father's day?

\texttt{Exercise\ starts}

\textbf{Exercise 1}

Make a comparison and an intersection cloud using bi-grams.

\textbf{Exercise 2}

Create a tri-gram word cloud.

\textbf{Exercise 3}

Print an igraph with more than 300 occurrences of bi-grams. Note the
differences and how this new graph gives you more insights.

\textbf{Exercise 4}

Analyze and visualize the frequency distribution of tri-grams across the
dataset.

\textbf{Exercise 5}

Take the iPad and Kindle subsets and make the word graphs for both of
them.

\texttt{Exercise\ ends}

\textbf{Conclusion} \emph{Conclusion starts}

This chapter was a walkthrough a text analysis from the importing of the basic data
and pre-processing, to more advanced concepts like the vector space model
and word frequencies. We looked at how the transformation of text into
numerical data enables us to extract information without the need to
read large volumes of text. The chapter highlighted the importance of steps like
normalization, stop words, punctuation, numbers, white space removal, and
vocabulary reduction in general. The application of
these techniques to word clouds, frequency plots, word networks,
and community detection provided a visual understanding of the data,
revealing underlying themes and patterns. We also plotted common words used for
different products as well as distinct words used to describe specific
products for easier identification. In this context, we examined the
importance of using n-grams for more nuanced analyses. We
explored how word frequencies can help develop SEO strategies, underscoring
the significance of selecting the right keywords to improve online
visibility and customer engagement.

\emph{Conclusion ends}
